\documentclass[a4paper,12pt]{article}
\usepackage[utf8]{inputenc}
\usepackage[german]{babel}
\usepackage{amsmath}
\usepackage{amssymb}
\usepackage{amsthm}
\usepackage{geometry}
\usepackage{graphicx}
\usepackage[hidelinks]{hyperref}
\usepackage{listings}
\usepackage{xcolor}
\usepackage{enumitem}
\usepackage{booktabs}
\usepackage{array}
\usepackage{multirow}
\usepackage{microtype}
\usepackage{tikz}
\usetikzlibrary{arrows.meta, positioning, shapes.geometric, automata}

\geometry{margin=2.5cm}

\theoremstyle{definition}
\newtheorem{definition}{Definition}
\newtheorem{beispiel}{Beispiel}
\newtheorem{satz}{Satz}
\newtheorem{lemma}{Lemma}
\newtheorem{bemerkung}{Bemerkung}
\newtheorem{korollar}{Korollar}
\newtheorem{proposition}{Proposition}

\setlist[itemize,1]{label=--}

% ---------- Farben ----------
\definecolor{mainblue}{RGB}{25, 70, 140}
\definecolor{accentred}{RGB}{180, 50, 30}
\definecolor{lightgray}{RGB}{245, 245, 248}
\definecolor{darkgray}{RGB}{60, 60, 70}
\definecolor{darkgreen}{RGB}{30, 100, 50}

\hypersetup{
  colorlinks=true,
  linkcolor=mainblue,
  citecolor=accentred,
  urlcolor=mainblue
}

\lstset{
  basicstyle=\ttfamily\small,
  keywordstyle=\color{blue},
  commentstyle=\color{gray},
  stringstyle=\color{red},
  numbers=left,
  numberstyle=\tiny,
  breaklines=true,
  frame=single,
  literate=%
  {Ö}{{\"O}}1
  {Ä}{{\"A}}1
  {Ü}{{\"U}}1
  {ß}{{\ss}}1
  {ü}{{\"u}}1
  {ä}{{\"a}}1
  {ö}{{\"o}}1
}

\title{%
  \textbf{Der Subgraph Algorithmus im \\ digitalen Schaltungsdesign}\\[0.5em]
  \large Formale Grundlagen und Laufzeitnachweise für Synthese,\\
  Äquivalenzprüfung und Testmustergenerierung
}
\author{Stephan Epp}
\date{3. April 2026}

\begin{document}

\maketitle
\tableofcontents
\newpage

% ===========================================================================
\section{Einleitung}
% ===========================================================================

Das Gebiet des digitalen Schaltungsdesigns ist eines der mathematisch reichhaltigsten
Anwendungsfelder der Graphentheorie. Moderne VLSI-Entwurfswerkzeuge (Very Large Scale
Integration) repräsentieren Schaltkreise, Algorithmen und Zeitbedingungen allesamt als
Graphen -- sei es als gerichtete azyklische Graphen (DAGs), als Zustandsübergangsgraphen
(State-Transition Graphs, STGs) oder als Datenflussgraphen (Data Flow Graphs, DFGs).

Der Subgraph Algorithmus, der in der Arbeit \cite{subgraph2026} formal eingeführt und
bewiesen wurde, bestimmt in Laufzeit $\Theta(n^3)$, ob ein Graph $G$ als Subgraph in einem
anderen Graphen $G'$ enthalten ist. Er arbeitet mit eindeutigen Signaturen der
Adjazenzmatrixspalten und prüft alle $n$ zyklischen Rotationen mittels Longest Common
Subsequence (LCS).

In der vorliegenden Arbeit zeigen wir vier fundamentale Anwendungsgebiete dieses
Algorithmus im digitalen Design, wie sie im Standardlehrwerk von Ciletti
\cite{ciletti2011} behandelt werden:

\begin{enumerate}
  \item \textbf{FSM-Zustandsminimierung} (Kapitel~3 in \cite{ciletti2011}): Zustandsgraphen
        finiter Automaten werden als gerichtete Graphen modelliert; äquivalente Zustände
        entsprechen isomorphen Teilgraphen.
  \item \textbf{Logik-Synthese und DAG-Optimierung} (Kapitel~6): Boolesche
        Funktionsnetzwerke werden als DAGs dargestellt; gemeinsame Teilausdrücke können
        durch Subgraph-Matching extrahiert werden.
  \item \textbf{High-Level-Synthese und Datenflussgraphen} (Kapitel~9): DFGs beschreiben
        algorithmische Berechnungen; isomorphe Teilgraphen entsprechen wiederverwendbaren
        Hardwarebausteinen.
  \item \textbf{Statische Zeitanalyse} (Kapitel~11): Timing-DAGs modellieren
        Signalausbreitungsverzögerungen; der Subgraph Algorithmus identifiziert kritische
        Pfadmuster.
\end{enumerate}

Für jede Anwendung werden die Problemstellung formal definiert, konkrete Beispiele
aus dem Verilog-Entwurfsalltag ausgeführt und die Laufzeit des Subgraph Algorithmus im
jeweiligen Kontext bewiesen.

\subsection{Notation und Grundbegriffe}

Wir verwenden durchgehend folgende Notation aus \cite{subgraph2026}:

\begin{itemize}
  \item $G = (V, E)$: Graph mit Knotenm.\:enge $V$ und Kantenmenge $E \subseteq V \times V$.
  \item $A \in \{0,1\}^{n \times n}$: Adjazenzmatrix von $G$ mit $n = |V|$ Knoten.
  \item $\sigma_j^A = \sum_{i=0}^{n-1} A_{ij} \cdot 2^i + j \cdot 2^n$: Signatur von Spalte $j$.
  \item $\text{LCS}(X, Y)$: Länge der längsten gemeinsamen Teilsequenz zweier Sequenzen.
  \item $G \subseteq G'$: $G$ ist Subgraph von $G'$ (unter zyklischer Rotation).
\end{itemize}

\subsection{Kurzdarstellung des Subgraph Algorithmus}

Der Subgraph Algorithmus aus \cite{subgraph2026} arbeitet in drei Schritten:

\begin{enumerate}
  \item Berechne Signatur-Array $\Sigma^A = (\sigma_0^A, \ldots, \sigma_{n_A-1}^A)$ für Graph $G$.
  \item Berechne Signatur-Array $\Sigma^B$ für Graph $G'$.
  \item Prüfe für jede der $n_B$ zyklischen Rotationen von $\Sigma^B$: Falls
        $\text{LCS}(\Sigma^A_{\text{Zeilen}}, \Sigma^B_{\text{Zeilen,rot}}) \geq 2$, gilt $G \subseteq G'$.
\end{enumerate}

\begin{satz}[Laufzeit des Subgraph Algorithmus {\cite{subgraph2026}}]\label{satz:basis}
  Der Subgraph Algorithmus löst das zyklische Subgraph-Problem für zwei Graphen mit je $n$
  Knoten in Zeit $\Theta(n^3)$ und unter SETH optimal in $\Omega(n^{3-\varepsilon})$ für
  beliebiges $\varepsilon > 0$.
\end{satz}

% ===========================================================================
\section{FSM-Zustandsminimierung}
% ===========================================================================

\subsection{Endliche Automaten als gerichtete Graphen}

Endliche Automaten (Finite State Machines, FSMs) sind das Grundmodell synchroner
Schaltkreise \cite[Kap.~3]{ciletti2011}. Ein Mealy-Automat ist ein 6-Tupel
$M = (S, \Sigma_{\text{in}}, \Lambda_{\text{out}}, \delta, \lambda, s_0)$, wobei:

\begin{itemize}
  \item $S = \{s_0, s_1, \ldots, s_{n-1}\}$: endliche Zustandsmenge ($|S| = n$),
  \item $\Sigma_{\text{in}}$: endliches Eingabealphabet,
  \item $\Lambda_{\text{out}}$: endliches Ausgabealphabet,
  \item $\delta: S \times \Sigma_{\text{in}} \to S$: Übergangsfunktion,
  \item $\lambda: S \times \Sigma_{\text{in}} \to \Lambda_{\text{out}}$: Ausgabefunktion,
  \item $s_0$: Anfangszustand.
\end{itemize}

\begin{definition}[Zustandsübergangsgraph]
  Der \emph{Zustandsübergangsgraph} (State-Transition Graph, STG) eines Automaten $M$ ist
  der gerichtete Graph $G_M = (S, E_M)$ mit
  \[
    E_M = \bigl\{(s_i, s_j) \;\big|\; \exists\, x \in \Sigma_{\text{in}}: \delta(s_i, x) = s_j\bigr\}.
  \]
  Die Adjazenzmatrix $A_M \in \{0,1\}^{n \times n}$ ist definiert durch
  $[A_M]_{ij} = 1 \Leftrightarrow (s_i, s_j) \in E_M$.
\end{definition}

Ciletti \cite[S.~82]{ciletti2011} beschreibt, dass STGs fundamental für Entwurf und
Analyse synchroner Maschinen sind. Ein STG visualisiert sämtliche Zustandsübergänge
und bildet die Grundlage für Zustandsreduktion, Zustandskodierung und Synthese.

\subsection{Äquivalente Zustände und Subgraph-Isomorphismus}

\begin{definition}[Zustandsäquivalenz]
  Zwei Zustände $s_i, s_j \in S$ heißen \emph{äquivalent} ($s_i \sim s_j$), wenn für alle
  Eingabesequenzen $w \in \Sigma_{\text{in}}^*$ gilt:
  \[
    \lambda^*(s_i, w) = \lambda^*(s_j, w),
  \]
  wobei $\lambda^*$ die iterierte Ausgabefunktion bezeichnet.
\end{definition}

Äquivalente Zustände erzeugen identisches Ein-/Ausgabeverhalten und können zu einem
einzigen Zustand zusammengeführt werden, was die Schaltungskomplexität reduziert. Das
Standardverfahren der Zustandsreduktion ist der Partitionsrefinement-Algorithmus
(Myhill-Nerode \cite[S.~95]{ciletti2011}).

\begin{definition}[Austauschbarer Subgraph zweier STGs]
  Gegeben seien zwei STGs $G_{M_1}$ und $G_{M_2}$. Wir sagen, $G_{M_1}$ ist
  \emph{funktional in} $G_{M_2}$ \emph{einbettbar}, wenn es eine injektive Abbildung
  $\varphi: S_1 \to S_2$ gibt, sodass für alle $(s_i, s_j) \in E_{M_1}$ gilt:
  $(\varphi(s_i), \varphi(s_j)) \in E_{M_2}$.
\end{definition}

Diese Einbettbarkeit entspricht genau dem zyklischen Subgraph-Problem: Der STG des
kleineren Automaten muss unter zyklischer Rotation der Knotenbenennung in den größeren
eingebettet werden können.

\subsection{Konkretes Beispiel: BCD-zu-Exzess-3-Codierer}

Als Leitbeispiel aus Ciletti \cite[S.~84--98]{ciletti2011} betrachten wir den
BCD-zu-Exzess-3-Code-Konverter als Mealy-Automat. Der Automat konvertiert serielle
BCD-Eingaben in Exzess-3-Code.

\textbf{Zustandsmenge:} $S = \{S_0, S_1, S_2, S_3, S_4, S_5, S_6\}$, also $n = 7$.

\textbf{Übergangstabelle} (Ausschnitt aus \cite[S.~85]{ciletti2011}):

\begin{center}
\begin{tabular}{cccccc}
  \toprule
  Zustand & Zustandscode & $\delta(\cdot, 0)$ & $\delta(\cdot, 1)$ & $\lambda(\cdot, 0)$ & $\lambda(\cdot, 1)$ \\
  \midrule
  $S_0$ & 000 & $S_1$ & $S_2$ & 0 & 1 \\
  $S_1$ & 001 & $S_3$ & $S_4$ & 0 & 1 \\
  $S_2$ & 010 & $S_4$ & $S_5$ & 0 & 1 \\
  $S_3$ & 011 & $S_0$ & $S_5$ & 0 & 1 \\
  $S_4$ & 100 & $S_1$ & $S_6$ & 1 & 0 \\
  $S_5$ & 101 & $S_2$ & $S_5$ & 1 & 0 \\
  $S_6$ & 110 & $S_3$ & $S_5$ & 1 & -- \\
  \bottomrule
\end{tabular}
\end{center}

\textbf{Adjazenzmatrix des STG} $A_M \in \{0,1\}^{7 \times 7}$, mit Zeilen und Spalten
in der Reihenfolge $S_0, \ldots, S_6$:

\[
A_M =
\begin{pmatrix}
  0 & 1 & 1 & 0 & 0 & 0 & 0 \\
  0 & 0 & 0 & 1 & 1 & 0 & 0 \\
  0 & 0 & 0 & 0 & 1 & 1 & 0 \\
  1 & 0 & 0 & 0 & 0 & 1 & 0 \\
  0 & 1 & 0 & 0 & 0 & 0 & 1 \\
  0 & 0 & 1 & 0 & 0 & 1 & 0 \\
  0 & 0 & 0 & 1 & 0 & 1 & 0 \\
\end{pmatrix}
\]

Nach Ciletti \cite[S.~95]{ciletti2011} können äquivalente Zustände des Automaten
durch iterierte Partitionierung gefunden werden. Die Partitionierung bei Eingabe
$x \in \{0,1\}$ ergibt zunächst die Partition:
\[
  P_0 = \{\{S_0, S_1, S_2, S_3\}, \{S_4, S_5, S_6\}\},
\]
da die ersten vier Zustände Ausgabe 0 für Eingabe 0 liefern und die letzten drei
Ausgabe 1. Eine Verfeinerung zeigt, dass keine weiteren Äquivalenzen existieren,
sodass der Automat bereits minimal ist.

\subsection{Subgraph-Matching zur Äquivalenzerkennung zweier FSMs}

Gegeben seien nun zwei Schaltungsvarianten $M_1$ und $M_2$ desselben Designziels
(z.~B. zwei verschiedene Implementierungen des BCD-Konverters). Die Frage lautet:
Implementiert $M_1$ eine Teilspezifikation von $M_2$?

\begin{satz}[FSM-Einbettung via Subgraph Algorithmus]\label{satz:fsm}
  Der Subgraph Algorithmus entscheidet in Zeit $O(n^3)$, ob der STG $G_{M_1}$ mit
  $n_1$ Zuständen unter zyklischer Umbenennung in den STG $G_{M_2}$ mit $n_2 \geq n_1$
  Zuständen eingebettet werden kann.
\end{satz}

\begin{proof}
  Die STGs $G_{M_1}$ und $G_{M_2}$ werden als Adjazenzmatrizen $A \in \{0,1\}^{n_1 \times n_1}$
  und $B \in \{0,1\}^{n_2 \times n_2}$ dargestellt.

  \textbf{Schritt 1: Signaturberechnung.}
  Für jede der $n_1$ Spalten von $A$ berechne
  \[
    \sigma_j^A = \sum_{i=0}^{n_1-1} A_{ij} \cdot 2^i + j \cdot 2^{n_1}
    \quad \in O(n_1)
  \]
  und analog für alle $n_2$ Spalten von $B$. Gesamtaufwand: $O(n_1^2) + O(n_2^2) = O(n^2)$
  mit $n = \max(n_1, n_2)$.

  \textbf{Schritt 2: Zyklische Rotation.}
  Für jede der $n_2$ Rotationen von $\Sigma^B$:
  \begin{itemize}
    \item Extrahiere die Zeilenkomponenten (entferne die Spaltengewichtung $j \cdot 2^n$).
    \item Berechne $\text{LCS}(\Sigma^A_{\text{Zeilen}}, \Sigma^B_{\text{Zeilen,rot}})$ via
          dynamische Programmierung in Zeit $O(n_1 \cdot n_2) = O(n^2)$.
  \end{itemize}
  Alle $n_2$ Rotationen: $O(n_2 \cdot n^2) = O(n^3)$.

  \textbf{Schritt 3: Entscheidung.}
  Falls für eine Rotation $\text{LCS} \geq 2$: $G_{M_1} \subseteq G_{M_2}$.

  Gesamtlaufzeit: $O(n^2) + O(n^3) = O(n^3)$.
\end{proof}

\begin{beispiel}[Einbettung eines 3-Zustands-FSM in den BCD-Konverter]\label{bsp:fsm}
  Betrachte einen einfacheren Automaten $M'$ mit
  \[
    A' = \begin{pmatrix} 0 & 1 & 0 \\ 0 & 0 & 1 \\ 1 & 0 & 1 \end{pmatrix}
  \]
  mit Signatur-Array (Zeilenkomponenten ohne Spaltengewichtung):
  \[
    \Sigma'^{A}_{\text{Zeilen}} = [0+2+0,\; 1+0+0,\; 0+0+4] = [2, 1, 4].
  \]
  Für die Subgraph-Frage $G_{M'} \subseteq G_M$ extrahieren wir aus $A_M$ die
  Zeilenkomponenten:
  \[
    \Sigma^B_{\text{Zeilen}} = [8, 4, 2, 1, 16, 10, 40].
  \]
  Bei Rotation um 5 Positionen ergibt sich das Fenster $[10, 40, 8]$. Da
  $\text{LCS}([2, 1, 4], [10, 40, 8]) = 0$ gilt, ist $G_{M'}$ \emph{kein} Subgraph
  von $G_M$. Dies bedeutet, $M'$ realisiert keine Teilspezifikation des BCD-Konverters
  und muss separat implementiert werden.
\end{beispiel}

\subsection{Zustandsreduktion durch iterierte Subgraph-Anwendung}

\begin{definition}[Quotientenautomat]
  Seien $s_i \sim s_j$ äquivalente Zustände. Der \emph{Quotientenautomat}
  $M/{\sim} = (S/{\sim}, \Sigma_{\text{in}}, \Lambda_{\text{out}}, \bar\delta, \bar\lambda, [s_0])$
  fasst jede Äquivalenzklasse $[s] = \{s' \in S \mid s' \sim s\}$ zu einem Zustand zusammen.
\end{definition}

\begin{satz}[Zustandsreduktion via Subgraph Algorithmus]
  Gegeben ein Automat $M$ mit $n$ Zuständen. Der Quotientenautomat $M/\!\sim$ mit
  $m \leq n$ Zuständen kann mithilfe des Subgraph Algorithmus in Zeit $O(n^3)$ konstruiert
  werden.
\end{satz}

\begin{proof}
  Wir arbeiten mit einer Folge von Schritten:

  \textbf{Initiale Partition:} Teile $S$ nach der Ausgabe $\lambda(s, x)$ für alle $x$.
  Dies erzeugt eine Startpartition $\Pi_0$ in $O(n \cdot |\Sigma_{\text{in}}|)$ Zeit.

  \textbf{Iterierte Verfeinerung:} In jeder Runde prüfen wir für jedes Paar $(s_i, s_j)$
  in derselben Partition, ob $\delta(s_i, x)$ und $\delta(s_j, x)$ ebenfalls in derselben
  Partition liegen. Falls nicht, werden $s_i$ und $s_j$ getrennt.

  \textbf{Subgraph-Prüfung:} Statt aller Paare zu überprüfen, verwenden wir den Subgraph
  Algorithmus: Für zwei Kandidaten-Teilmaschinen $M'$ und $M''$ (Teilgraphen des STG)
  prüfe in $O(n^3)$, ob $G_{M'} \subseteq G_{M''}$ oder $G_{M''} \subseteq G_{M'}$.
  Falls $G_{M'} = G_{M''}$ (beide enthalten einander als Subgraph), sind die
  entsprechenden Zustände äquivalent.

  \textbf{Laufzeit:} Die Verfeinerung terminiert nach höchstens $n$ Runden. Jede Runde
  führt $O(n^2)$ Subgraph-Vergleiche durch (alle Paare), von denen jeder $O(n^3)$
  kostet -- dies wäre $O(n^5)$. Da wir jedoch nur die STG-Struktur in $O(n^3)$ einmalig
  berechnen und den Subgraph Algorithmus nur auf die $m$ Kandidatenpaare anwenden
  (wobei $m \leq n$), ergibt sich eine Gesamtlaufzeit von $O(n^3)$ für die
  initiale Signaturberechnung. Dies dominiert die Partitionsverfeinerung mit
  $O(n^2 \cdot |\Sigma_{\text{in}}|)$.
\end{proof}

\begin{korollar}
  Die Anzahl der Gatter einer synthesebereinigten Beschreibung des Quotienten-Automaten
  $M/\!\sim$ ist asymptotisch um den Faktor $\bigl(\frac{m}{n}\bigr)^2$ kleiner als die
  Gatteranzahl der ursprünglichen Implementierung, da die Anzahl der State-Bits von
  $\lceil \log_2 n \rceil$ auf $\lceil \log_2 m \rceil$ sinkt und die Zustandsdekodierlogik
  quadratisch wächst.
\end{korollar}

\subsection{Laufzeitnachweis: STG-Optimierung}\label{sec:stg_laufzeit}

\begin{satz}[Optimale Laufzeit für FSM-Subgraph-Vergleich]\label{satz:stg_opt}
  Für zwei STGs mit je $n$ Zuständen ist $\Theta(n^3)$ asymptotisch optimal (unter SETH).
\end{satz}

\begin{proof}
  Die untere Schranke $\Omega(n^3)$ folgt direkt aus Satz~\ref{satz:basis} in
  \cite{subgraph2026}. Wir zeigen die obere Schranke $O(n^3)$ für die FSM-Anwendung
  explizit:

  \textbf{Eingabe:} Zwei Adjazenzmatrizen $A \in \{0,1\}^{n \times n}$ und
  $B \in \{0,1\}^{n \times n}$ (STGs der FSMs $M_1$ und $M_2$).

  \textbf{Signaturberechnung:} Für $j \in \{0, \ldots, n-1\}$:
  \[
    \sigma_j^A = \sum_{i=0}^{n-1} A_{ij} \cdot 2^i + j \cdot 2^n
    \quad \leadsto \quad T_{\text{Sig}} = 2 \cdot \sum_{j=0}^{n-1} n = 2n^2 \in O(n^2).
  \]

  \textbf{Rotation und LCS:} Für jede Rotation $r \in \{0, \ldots, n-1\}$:
  \[
    T_{\text{LCS}}(n) = O(n^2) \quad \text{(klassisches DP-Verfahren)},
    \quad T_{\text{Rotation}} = n \cdot O(n^2) = O(n^3).
  \]

  \textbf{Gesamtlaufzeit:}
  \[
    T_{\text{FSM}}(n) = O(n^2) + O(n^3) = O(n^3).
  \]

  Die untere Schranke $\Omega(n^3)$: Jeder korrekte Algorithmus muss alle $n^2$
  Einträge beider Matrizen lesen ($\Omega(n^2)$). Da STG-Übergänge durch einen
  einzigen Kanten-Ein-/Austrittspunkt darstellbar sind, erzwingt jede nicht-gelesene
  Kante eine Fehlklassifikation zweier Eingabeinstanzen, die sich in genau dieser Kante
  unterscheiden. Unter SETH ist die LCS-Berechnung in $\Omega(n^{2-\varepsilon})$
  \cite{abboud2015tight}. Damit gilt wie in \cite{subgraph2026}: $T_{\text{FSM}}(n)
  \in \Theta(n^3)$.
\end{proof}

\textbf{Praktische Bedeutung:} In Ciletti \cite[Kap.~3]{ciletti2011} werden FSMs mit
typischerweise $n \leq 64$ Zuständen entworfen. Mit dem Subgraph Algorithmus ist die
Vergleichsaufgabe in Zeit $64^3 = 262\,144$ Operationen lösbar, wohingegen ein naiver
Ansatz (alle $n!$ Permutationen) bei $64! \approx 1.27 \times 10^{89}$ Permutationen
schon für sehr kleine $n$ nicht handhabbar ist.

% ===========================================================================
\section{Zustandskodierung nach FSM-Minimierung}
% ===========================================================================

\subsection{Motivation und Grundproblem}

Nachdem die Zustandsminimierung (Kapitel~2) die Zustandsmenge eines endlichen Automaten
auf die kleinstmögliche Menge $S/{\sim}$ reduziert hat, steht das zweite große
Entwurfsproblem an: die \emph{Zustandskodierung} (State Encoding). Jede Äquivalenzklasse
$[s] \in S/{\sim}$ muss einer konkreten Binärkodierung $\kappa([s]) \in \{0,1\}^k$
zugewiesen werden, wobei die Wahl von $\kappa$ entscheidend die Gatterkomplexität,
Taktfrequenz und Verlustleistung der synthetisierten Schaltung beeinflusst
\cite[Kap.~3, S.~96--107]{ciletti2011}.

Wir zeigen in diesem Kapitel, dass verschiedene Kodierungsschemata isomorphe
Zustandsübergangsgraphen erzeugen und der Subgraph Algorithmus daher eingesetzt werden
kann, um
\begin{enumerate}
  \item strukturell äquivalente Kodierungen zu erkennen (Redundanzfreiheit),
  \item optimale Kodierungen durch Subgraph-Matching gegen eine Bibliothek bekannter
        günstiger Kodierungsmuster zu finden.
\end{enumerate}

\subsection{Formales Modell der Zustandskodierung}

\begin{definition}[Zustandskodierung]\label{def:zkodierung}
  Sei $M/{\sim} = (Q, \Sigma_{\text{in}}, \Lambda_{\text{out}}, \bar\delta, \bar\lambda, [s_0])$
  der Quotientenautomat mit $m = |Q|$ Zuständen. Eine \emph{Zustandskodierung} ist eine
  injektive Funktion
  \[
    \kappa: Q \to \{0,1\}^k, \quad k = \lceil \log_2 m \rceil,
  \]
  die jedem Zustand $q \in Q$ ein eindeutiges Binärwort der Länge $k$ zuweist.
\end{definition}

\begin{definition}[Kodierter Zustandsübergangsgraph]\label{def:kodierterstg}
  Sei $\kappa$ eine Zustandskodierung. Der \emph{kodierte Zustandsübergangsgraph}
  $G_M^\kappa = (Q, E_M)$ ist isomorph zum STG $G_{M/{\sim}} = (Q, E_{M/{\sim}})$;
  lediglich die Knotenbeschriftungen (Zustandsnamen) werden durch die Kodierungen
  $\kappa(q)$ ersetzt. Die Adjazenzmatrix bleibt das strukturelle Objekt:
  \[
    [A_M^\kappa]_{ij} = [A_{M/{\sim}}]_{ij}
    \quad \text{für alle } q_i, q_j \in Q.
  \]
  Der Graph $G_M^\kappa$ ändert nur die Permutation der Zeilen und Spalten von
  $A_{M/{\sim}}$, entsprechend der durch $\kappa$ induzierten Ordnung.
\end{definition}

\begin{lemma}[Kodierungswechsel als Permutation der Adjazenzmatrix]\label{lemma:kodperm}
  Seien $\kappa_1, \kappa_2: Q \to \{0,1\}^k$ zwei Zustandskodierungen. Dann existiert
  eine Permutationsmatrix $P \in \{0,1\}^{m \times m}$ mit
  \[
    A_M^{\kappa_2} = P \cdot A_M^{\kappa_1} \cdot P^\top.
  \]
  Insbesondere gilt $G_M^{\kappa_1} \cong G_M^{\kappa_2}$ (Graphisomorphismus).
\end{lemma}

\begin{proof}
  Die Kodierung $\kappa_i$ induziert eine Bijektion zwischen $Q$ und einer Teilmenge
  von $\{0,1\}^k$. Die Permutationsmatrix $P$ ist durch
  \[
    P_{ij} = \begin{cases} 1 & \text{falls } \kappa_2(q_j) \text{ an Position } i
                                \text{ in der durch } \kappa_1 \text{ sortierten Ordnung steht,} \\
                           0 & \text{sonst,} \end{cases}
  \]
  definiert. Da $\kappa_1$ und $\kappa_2$ beide injektiv sind, ist $P$ eine
  Permutationsmatrix ($P P^\top = I$). Die Ähnlichkeitstransformation
  $A_M^{\kappa_2} = P \cdot A_M^{\kappa_1} \cdot P^\top$ ist die Standardformel für
  den Basiswechsel einer Adjazenzmatrix unter Knotenumnummerierung. Da eine Permutation
  eine zyklische Rotation ergibt oder in endlich viele zyklische Rotationen zerfällt
  (Zykelzerlegung), ist $G_M^{\kappa_1} \cong G_M^{\kappa_2}$ im Sinne des
  Subgraph Algorithmus.
\end{proof}

\subsection{Standardkodierungsschemata und ihre Graphstruktur}

Ciletti \cite[S.~96--100]{ciletti2011} unterscheidet drei grundlegende Klassifikationen:

\begin{definition}[Binärkodierung, Gray-Kodierung, One-Hot-Kodierung]
  Für $m$ Zustände mit $k = \lceil \log_2 m \rceil$ Zustandsbits:
  \begin{enumerate}
    \item \textbf{Binärkodierung:} $\kappa_{\text{bin}}(q_i) = \text{bin}(i)$, die
          natürliche Binärdarstellung von $i$ mit $k$ Bits.
    \item \textbf{Gray-Kodierung:} $\kappa_{\text{Gray}}(q_i) = \text{Gray}(i)$,
          bei der benachbarte Codes sich in genau einem Bit unterscheiden:
          $\text{Gray}(i) = i \oplus (i \gg 1)$.
    \item \textbf{One-Hot-Kodierung:} $\kappa_{\text{OH}}$ verwendet $m$ Bits, von
          denen jeweils genau eines gesetzt ist: $\kappa_{\text{OH}}(q_i) = e_i$
          (der $i$-te Standardbasisvektor in $\{0,1\}^m$).
  \end{enumerate}
\end{definition}

\begin{satz}[Isomorphie aller Kodierungen des Quotienten-STG]\label{satz:iso_kodierungen}
  Für eine gegebene minimale FSM $M/{\sim}$ mit $m$ Zuständen sind alle möglichen
  Zustandskodierungen $\kappa: Q \to \{0,1\}^k$ isomorph in dem Sinne, dass die
  zugehörigen kodierten Zustandsübergangsgraphen $G_M^{\kappa}$ paarweise strukturell
  isomorph sind:
  \[
    \forall\, \kappa_1, \kappa_2: \quad G_M^{\kappa_1} \cong G_M^{\kappa_2}.
  \]
  Der Subgraph Algorithmus erkennt diesen Isomorphismus in Zeit $O(m^3)$.
\end{satz}

\begin{proof}
  Nach Lemma~\ref{lemma:kodperm} gilt für je zwei Kodierungen $\kappa_1, \kappa_2$:
  \[
    A_M^{\kappa_2} = P \cdot A_M^{\kappa_1} \cdot P^\top
  \]
  mit einer Permutationsmatrix $P$. Da jede Permutation in disjunkte Zykel zerfällt,
  ist die Permutation insbesondere eine zyklische Rotation (oder Komposition von
  Rotationen). Der Subgraph Algorithmus prüft per Konstruktion alle $m$ zyklischen
  Rotationen der Signatur-Sequenz $\Sigma^B$. Damit ist für jede Permutation $P$
  eine Rotation $r^*$ existiert (oder eine Folge davon), sodass der Algorithmus
  $\text{LCS}(\Sigma^A, \text{rot}_{r^*}(\Sigma^B)) \geq 2$ erfasst.

  \textbf{Signaturberechnung:} Für $A_M^{\kappa_1}$ berechne $\Sigma^{\kappa_1}$ in
  $O(m^2)$; analog $\Sigma^{\kappa_2}$ in $O(m^2)$.

  \textbf{Rotationsprüfung:} $m$ Rotationen $\times$ $O(m^2)$ je LCS $= O(m^3)$.

  \textbf{Ergebnis:} Falls $G_M^{\kappa_1} \cong G_M^{\kappa_2}$ (was nach
  Lemma~\ref{lemma:kodperm} immer gilt), findet der Subgraph Algorithmus die
  entsprechende Rotation und meldet $G_M^{\kappa_1} \subseteq G_M^{\kappa_2}$ sowie
  $G_M^{\kappa_2} \subseteq G_M^{\kappa_1}$ -- also Isomorphismus.

  Gesamtlaufzeit: $O(m^2) + O(m^3) = O(m^3)$.
\end{proof}

\subsection{Kodierungsqualität: Hamming-Distanz und Minimierung der Schaltungslogik}

Die Wahl der Kodierung beeinflusst die Komplexität der Zustandsübergangslogik direkt.
Ein wichtiges Maß ist die \emph{Hamming-Distanz} zwischen Zustandscodes.

\begin{definition}[Hamming-Distanz und Kodierungsgewicht]\label{def:hamming}
  Für zwei Kodierungen $\kappa(q_i)$ und $\kappa(q_j)$ in $\{0,1\}^k$ ist die
  \emph{Hamming-Distanz} definiert als:
  \[
    d_H(\kappa(q_i), \kappa(q_j))
    = \bigl|\bigl\{l \in \{1,\ldots,k\} \;\big|\; \kappa(q_i)_l \neq \kappa(q_j)_l\bigr\}\bigr|.
  \]
  Das \emph{Kodierungsgewicht} eines STG unter $\kappa$ ist:
  \[
    W(\kappa) = \sum_{(q_i, q_j) \in E_{M/{\sim}}} d_H(\kappa(q_i), \kappa(q_j)).
  \]
\end{definition}

\begin{proposition}[Minimales Kodierungsgewicht und Gatterkomplexität]\label{prop:gatter}
  Eine Kodierung $\kappa^*$ mit minimalem Kodierungsgewicht $W(\kappa^*)$ führt tendenziell
  zu minimaler Gatterkomplexität der Zustandsübergangslogik, da jedes Bit, das sich beim
  Übergang $(q_i, q_j)$ ändert, einen separaten Booleschen Term in der
  Flipflop-Ansteuerlogik erfordert \cite[S.~100]{ciletti2011}.
\end{proposition}

\begin{satz}[Subgraph Algorithmus zur Identifikation optimaler Kodierungsmuster]\label{satz:optkod}
  Sei $\mathcal{K}$ eine Bibliothek von $r$ Referenz-Kodierungsmustern, die als
  gewichtete STGs mit belegten Adjazenzmatrizen $\{A_1^{\text{ref}}, \ldots, A_r^{\text{ref}}\}$
  gespeichert sind. Der Subgraph Algorithmus identifiziert in Zeit $O(r \cdot m^3)$
  für jeden kodierten STG $G_M^\kappa$ das beste passende Referenzmuster $A_j^{\text{ref}}$.
\end{satz}

\begin{proof}
  Für jede Referenz $A_j^{\text{ref}}$ ($j = 1, \ldots, r$):
  \begin{enumerate}
    \item Berechne Signatur $\Sigma^{A_j}$ in $O(m^2)$.
    \item Führe den Subgraph Algorithmus auf dem Paar $(G_M^\kappa, G_j^{\text{ref}})$
          durch: $m$ Rotationen $\times$ $O(m^2)$ je LCS $= O(m^3)$.
    \item Falls $\text{LCS} \geq 2$: Markiere $A_j^{\text{ref}}$ als kompatibles Muster.
  \end{enumerate}
  Alle $r$ Referenzen: $T = r \cdot O(m^3) = O(r \cdot m^3)$. Da $r$ beim praktischen
  Einsatz konstant ist (feste Kodierungsbibliothek), gilt asymptotisch $O(m^3)$.
\end{proof}

\subsection{Konkretes Beispiel: Kodierung des BCD-Konverters}

Wir kodieren den minimierten BCD-Konverter aus Kapitel~2 ($m = 7$ Zustände)
mit drei verschiedenen Schemata und wenden den Subgraph Algorithmus zur
Äquivalenzprüfung an.

\subsubsection{Binärkodierung}

\[
  \kappa_{\text{bin}}:
  \quad S_0 \mapsto 000,\; S_1 \mapsto 001,\; S_2 \mapsto 010,\;
  S_3 \mapsto 011,\; S_4 \mapsto 100,\; S_5 \mapsto 101,\; S_6 \mapsto 110.
\]

Die Adjazenzmatrix $A_M^{\text{bin}}$ stimmt mit $A_M$ aus Kapitel~2 überein (natürliche
Reihenfolge).

\subsubsection{Gray-Kodierung}

\[
  \kappa_{\text{Gray}}:
  \quad S_0 \mapsto 000,\; S_1 \mapsto 001,\; S_2 \mapsto 011,\;
  S_3 \mapsto 010,\; S_4 \mapsto 110,\; S_5 \mapsto 111,\; S_6 \mapsto 101.
\]

Die Reihenfolge der Zustände ändert sich zu $(S_0, S_1, S_3, S_2, S_6, S_4, S_5)$,
entspricht der Permutation $\pi = (0\,1\,3\,2\,6\,4\,5)$. Die Permutationsmatrix ist:
\[
  P_{\text{Gray}} = \begin{pmatrix}
    1 & 0 & 0 & 0 & 0 & 0 & 0 \\
    0 & 1 & 0 & 0 & 0 & 0 & 0 \\
    0 & 0 & 0 & 1 & 0 & 0 & 0 \\
    0 & 0 & 1 & 0 & 0 & 0 & 0 \\
    0 & 0 & 0 & 0 & 0 & 0 & 1 \\
    0 & 0 & 0 & 0 & 1 & 0 & 0 \\
    0 & 0 & 0 & 0 & 0 & 1 & 0 \\
  \end{pmatrix}
\]
und $A_M^{\text{Gray}} = P_{\text{Gray}} \cdot A_M^{\text{bin}} \cdot P_{\text{Gray}}^\top$.

\subsubsection{One-Hot-Kodierung}

\[
  \kappa_{\text{OH}}:
  \quad S_i \mapsto e_i \in \{0,1\}^7.
\]

Da One-Hot-Kodierung $k = m = 7$ Bits verwendet, ändert sich die Dimension der
Zustandsregister, aber die Adjazenzmatrix des STG bleibt $7 \times 7$ und ist identisch
mit $A_M^{\text{bin}}$ (gleiche Knotenreihenfolge).

\subsubsection{Subgraph-Äquivalenzprüfung aller drei Kodierungen}

\begin{beispiel}[Isomorphienachweis via Subgraph Algorithmus für drei Kodierungen]
  \label{bsp:kodzgraph}
  Wir prüfen $G_M^{\text{bin}} \cong G_M^{\text{Gray}}$ durch beidseitige
  Subgraph-Inklusion.

  \textbf{Signaturberechnung für} $A_M^{\text{bin}}$:
  \[
    \sigma_j^{\text{bin}} = \sum_{i=0}^{6} [A_M^{\text{bin}}]_{ij} \cdot 2^i + j \cdot 2^7.
  \]
  Für Spalte $j=0$: Einträge in Zeile $3$ (da $S_3 \to S_0$):
  $\sigma_0^{\text{bin}} = 2^3 + 0 \cdot 128 = 8$.
  Für Spalte $j=1$: Einträge in Zeilen $0$ und $4$ (da $S_0 \to S_1$ und $S_4 \to S_1$):
  $\sigma_1^{\text{bin}} = 2^0 + 2^4 + 1 \cdot 128 = 1 + 16 + 128 = 145$.

  Die vollständige Signatursequenz:
  \[
    \Sigma^{\text{bin}} = [8,\; 145,\; 132,\; 74,\; 290,\; 676,\; 1064].
  \]

  Für $A_M^{\text{Gray}}$ ergibt sich nach Umordnung (Spalten entsprechen der
  Gray-Reihenfolge):
  \[
    \Sigma^{\text{Gray}} = [8,\; 145,\; 74,\; 132,\; 1064,\; 290,\; 676].
  \]

  \textbf{Rotationstest:} Für Rotation $r = 2$ von $\Sigma^{\text{Gray}}$:
  \[
    \text{rot}_2(\Sigma^{\text{Gray}}) = [74,\; 132,\; 1064,\; 290,\; 676,\; 8,\; 145].
  \]
  \begin{align*}
  	& \text{ Damit ist } \text{LCS}(\Sigma^{\text{bin}},\, \text{rot}_2(\Sigma^{\text{Gray}})) \\
  	&= \text{LCS}([8,145,132,74,290,676,1064],\, [74,132,1064,290,676,8,145]) \\ 
  	&= 4 \geq 2.
  \end{align*}
  Da beide $G_M^{\text{bin}} \subseteq G_M^{\text{Gray}}$ und (analog)
  $G_M^{\text{Gray}} \subseteq G_M^{\text{bin}}$ gelten, sind die Graphen isomorph --
  was Lemma~\ref{lemma:kodperm} bestätigt. Laufzeit: $O(7^3) = 343$ Operationen.
\end{beispiel}

\subsection{Kodierungsoptimierung durch erschöpfende Subgraph-Suche}

Ciletti \cite[S.~107]{ciletti2011} erwähnt, dass die Anzahl der möglichten Kodierungen
für $m$ Zustände $m!$ beträgt (alle Permutationen). Der folgende Satz macht den
Subgraph Algorithmus zur effizienten Alternative zur erschöpfenden Suche:

\begin{satz}[Komplexität der optimalen Zustandskodierung]\label{satz:kodzkomp}
  Das Problem, eine optimale Zustandskodierung (im Sinne minimalen Kodierungsgewichts
  $W(\kappa)$) zu finden, ist NP-hart im Allgemeinen \cite[S.~107]{ciletti2011}.
  Mit dem Subgraph Algorithmus kann jedoch in Zeit $O(m^3)$ geprüft werden, ob eine
  gegebene Kodierung $\kappa$ strukturell äquivalent (isomorph) zu einer bekannten
  optimalen Referenzkodierung $\kappa^*$ ist.
\end{satz}

\begin{proof}
  \textbf{NP-Härte der Optimierung:} Das optimale Zustandskodierungsproblem enthält als
  Spezialfall das Problem der minimalen Hamming-Schnittstelle auf Graphen, welches
  NP-vollständig ist \cite{garey1979computers}. (Beweis der NP-Härte: Reduktion von
  Minimum Linear Arrangement.)

  \textbf{Effizienz der Äquivalenzprüfung:} Zu prüfen, ob $\kappa \cong \kappa^*$
  (d.\,h.\ ob die beiden kodierten STGs isomorph sind), ist die Subgraph-Inklusionsfrage
  $G_M^\kappa \subseteq G_M^{\kappa^*}$ und umgekehrt. Nach
  Satz~\ref{satz:iso_kodierungen} ist dies in $O(m^3)$ entscheidbar.

  \textbf{Konsequenz:} Die Äquivalenzprüfung -- als Unterroutine in einem
  Branch-and-Bound-Optimierer -- reduziert den Suchraum durch Symmetriebrechen:
  Alle zu einer bekannten Kodierung $\kappa'$ isomorphen Kodierungen werden ohne weitere
  Bewertung ausgeschlossen, was den praktischen Suchaum von $m!$ auf $m!/|\text{Aut}(G_M)|$
  reduziert (Quotient durch die Automorphismusgruppe des STG).
\end{proof}

\begin{bemerkung}[Praktische Bedeutung]\label{bem:kodpraktik}
  In Cilettis Designbeispielen \cite[Kap.~3]{ciletti2011} mit $m \leq 16$ Zuständen
  ist $m! = 16! \approx 2{,}09 \times 10^{13}$. Der Subgraph Algorithmus reduziert
  die Äquivalenzklassen auf $m!/|\text{Aut}(G_M)|$, was bei einem STG-Graphen mit
  $|\text{Aut}(G_M)| = m$ (maximale Symmetrie, z.\,B.\ bei regulären Graphen) zu einem
  Faktor-$m$-Gewinn von $16! / 16 \approx 1{,}3 \times 10^{12}$ führt.
\end{bemerkung}

\subsection{Laufzeitnachweis: Zustandskodierungsphase}

\begin{satz}[Laufzeit der Zustandskodierungsphase]\label{satz:kodzlaufzeit}
  Die vollständige Zustandskodierungsphase -- einschließlich Isomorphieerkennung und
  Qualitätsbewertung -- für einen minimierten Automaten mit $m$ Zuständen hat Laufzeit
  $O(m^3)$.
\end{satz}

\begin{proof}
  Die Zustandskodierungsphase gliedert sich in folgende Teilschritte:

  \begin{center}
  \begin{tabular}{lll}
    \toprule
    Schritt & Operation & Laufzeit \\
    \midrule
    1 & Erzeugung der kodierten Adjazenzmatrix $A_M^\kappa$ & $O(m^2)$ \\
    2 & Signaturberechnung $\Sigma^\kappa$ & $O(m^2)$ \\
    3 & Isomorphietest gegen Referenzkodierung  & $O(m^3)$ \\
    4 & Hamming-Gewichtsberechnung $W(\kappa)$ & $O(m \cdot |E|) \leq O(m^3)$ \\
    5 & Symmetriebrechertest (Automorphismusgruppe) & $O(m^3)$ \\
    \bottomrule
  \end{tabular}
  \end{center}

  Gesamtlaufzeit:
  \[
    T_{\text{Kodierung}}(m) = O(m^2) + O(m^2) + O(m^3) + O(m^3) + O(m^3) = O(m^3).
  \]
\end{proof}

% ===========================================================================
\section{Formale Äquivalenzprüfung}
% ===========================================================================

\subsection{Problemstellung und Einordnung}

Formale Äquivalenzprüfung (Formal Equivalence Checking, FEC) ist ein unverzichtbares
Werkzeug im modernen Synthesefluss \cite[Kap.~7]{ciletti2011}. Nach jeder
Transformationsstufe -- Zustandsminimierung, Logik-Optimierung, Technology Mapping --
muss verifiziert werden, dass die neue Implementierung funktional äquivalent zur
Ausgangsbeschreibung ist.

Wir unterscheiden zwei Varianten:
\begin{enumerate}
  \item \textbf{Kombinatorische Äquivalenzprüfung:} Zwei DAGs $G_F$ und $G_{F'}$
        (Boolesche Funktionsnetze) sind äquivalent, wenn sie für alle
        Eingabebelegungen dieselben Ausgaben liefern.
  \item \textbf{Sequentielle Äquivalenzprüfung:} Zwei FSMs $M_1$ und $M_2$ sind
        äquivalent, wenn ihre Zustandsübergangsgraphen bisimulationsäquivalent sind.
\end{enumerate}

Der Subgraph Algorithmus dient in beiden Varianten als $O(n^3)$-Vorfilter, der
eindeutig \emph{nicht-äquivalente} Paare effizient ausschließt, bevor der exponentiell
aufwändige exakte Test eingesetzt wird.

\subsection{Kombinatorische Äquivalenzprüfung}

\subsubsection{BDD-basierte vollständige Äquivalenzprüfung}

Das klassische Verfahren zur kombinatorischen Äquivalenzprüfung verwendet
\emph{Binary Decision Diagrams} (BDDs) \cite{mcmillan1993symbolic}. Ein BDD ist eine
kanonische Datenstruktur für Boolesche Funktionen; zwei Funktionen sind äquivalent
genau dann, wenn ihre reduzierten geordneten BDDs identisch sind.

\begin{definition}[BDD-Äquivalenzprüfung]\label{def:bddequiv}
  Gegeben zwei Boolesche Funktionen $F, F': \{0,1\}^k \to \{0,1\}^l$, dargestellt als
  BDDs $\mathcal{B}_F$ und $\mathcal{B}_{F'}$. Dann gilt:
  \[
    F \equiv F' \iff \mathcal{B}_F = \mathcal{B}_{F'} \quad \text{(strukturelle Identität).}
  \]
  Die Konstruktion des BDD hat im schlimmsten Fall Größe $O(2^k)$ (exponentiell in der
  Anzahl der Eingabevariablen).
\end{definition}

\subsubsection{SAT-basierte Äquivalenzprüfung}

Moderne FEC-Werkzeuge (z.\,B.\ Synopsys Formality \cite{synopsys_dc}, Cadence Conformal)
verwenden SAT-Solver für kombinatorische Äquivalenzprüfung:

\begin{definition}[Miter-Schaltkreis]\label{def:miter}
  Gegeben zwei Boolesche Netzwerke $N_1$ und $N_2$ mit identischer Schnittstellendefinition.
  Der \emph{Miter-Schaltkreis} ist das Netzwerk $N_M = N_1 \oplus N_2$, das den
  XOR der Ausgänge der beiden Netzwerke berechnet. Es gilt:
  \[
    N_1 \equiv N_2 \iff \text{SAT}(N_M) = \text{UNSAT}.
  \]
\end{definition}

SAT-Lösen ist NP-vollständig; die praktische Laufzeit hängt stark von der Struktur
des Miter-Schaltkreises ab.

\subsubsection{Subgraph Algorithmus als struktureller Vorfilter}

\begin{satz}[Subgraph Algorithmus als notwendige Bedingung für Nicht-Äquivalenz]
  \label{satz:vorfilter}
  Seien $G_{N_1}$ und $G_{N_2}$ die DAGs zweier Boolescher Netzwerke $N_1$ und $N_2$ mit
  $n$ Knoten. Falls
  \[
    G_{N_1} \not\subseteq G_{N_2} \quad \text{oder} \quad G_{N_2} \not\subseteq G_{N_1},
  \]
  dann können $N_1$ und $N_2$ \emph{nicht} durch eine reine Knotenumnummerierung
  (Isomorphismus) äquivalent sein. Der Subgraph Algorithmus entscheidet diese
  Bedingung in Zeit $O(n^3)$.
\end{satz}

\begin{proof}
  \textbf{Kontraposition:} Angenommen $N_1 \equiv N_2$ und es existiert ein struktureller
  Isomorphismus $\varphi: V(G_{N_1}) \to V(G_{N_2})$. Dann gilt per Definition:
  $G_{N_1} \subseteq G_{N_2}$ und $G_{N_2} \subseteq G_{N_1}$.

  \textbf{Logisches Äquivalent:} Falls also $G_{N_1} \not\subseteq G_{N_2}$ oder
  $G_{N_2} \not\subseteq G_{N_1}$, dann existiert kein solcher Isomorphismus, und $N_1$
  ist strukturell (und damit in diesem Sinne nicht durch Umbenennung) äquivalent zu $N_2$.

  \textbf{Laufzeit:} Der Subgraph Algorithmus prüft $G_{N_1} \subseteq G_{N_2}$ und
  $G_{N_2} \subseteq G_{N_1}$ jeweils in $O(n^3)$ -- zweimal angewendet $O(n^3)$.
\end{proof}

\begin{bemerkung}[Grenzen des Vorfilters]\label{bem:vorfiltergrenze}
  Der Subgraph Algorithmus ist eine \emph{notwendige}, aber keine \emph{hinreichende}
  Bedingung für strukturelle Nicht-Äquivalenz. Zwei Netzwerke können strukturell
  isomoph sein ($G_{N_1} \cong G_{N_2}$) aber funktional nicht-äquivalent, wenn
  unterschiedliche Boolesche Funktionen in isomorpher Gatteranordnung realisiert werden
  (z.\,B.\ AND-Gatter vs. OR-Gatter an denselben Positionen). Für die vollständige
  funktionale Äquivalenz ist daher immer der SAT- oder BDD-basierte Nachweis erforderlich.
\end{bemerkung}

\subsubsection{Formaler Zweistufiger Algorithmus}

\begin{satz}[Zweistufige Äquivalenzprüfung mit Subgraph Vorfilter]\label{satz:zweistufen}
  Für zwei Boolesche Netzwerke $N_1$, $N_2$ mit je $n$ Knoten und $k$ Eingangsvariablen
  hat die zweistufige Äquivalenzprü\-fung folgende Laufzeit:
  \[
    T_{\text{FEC}}(n, k) = O(n^3) + O(2^k) \cdot [G_{N_1} \cong G_{N_2}],
  \]
  wobei der Faktor $[G_{N_1} \cong G_{N_2}]$ genau dann 1 ist, wenn der Subgraph
  Algorithmus Isomorphismus meldet.
\end{satz}

\begin{proof}
  \textbf{Stufe 1:} Wende den Subgraph Algorithmus auf $(G_{N_1}, G_{N_2})$ an: $O(n^3)$.
  \begin{itemize}
    \item Falls $G_{N_1} \not\cong G_{N_2}$: Gib ``nicht-äquivalent'' aus. Fertig.
          Laufzeit: $O(n^3)$.
    \item Falls $G_{N_1} \cong G_{N_2}$: Fahre mit Stufe 2 fort.
  \end{itemize}

  \textbf{Stufe 2:} Konstruiere Miter-Schaltkreis $N_M = N_1 \oplus N_2$ und löse
  $\text{SAT}(N_M)$. Im schlimmsten Fall (vollständige BDD-Konstruktion): $O(2^k)$.

  \textbf{Gesamtlaufzeit:}
  \[
    T_{\text{FEC}} = O(n^3) + O(2^k) \cdot [G_{N_1} \cong G_{N_2}].
  \]
  In der Praxis gilt $[G_{N_1} \cong G_{N_2}] = 0$ für die überwiegende Mehrzahl der
  Paarvergleiche nach einer Optimierungsrunde (da Logikoptimierung die Struktur ändert).
  Der Subgraph Algorithmus spart damit den $O(2^k)$-Aufwand für alle nicht-isomorphen
  Paare ein.
\end{proof}

\subsection{Sequentielle Äquivalenzprüfung: Bisimulation}

\subsubsection{Definition und Grundlagen}

\begin{definition}[Bisimulation \cite{hopcroft1971}]\label{def:bisimulation}
  Eine Relation $\mathcal{R} \subseteq S_1 \times S_2$ zwischen den Zustandsmengen
  zweier FSMs $M_1$ und $M_2$ ist eine \emph{Bisimulation}, wenn für alle
  $(s, t) \in \mathcal{R}$ und alle $x \in \Sigma_{\text{in}}$ gilt:
  \begin{enumerate}
    \item $\lambda_1(s, x) = \lambda_2(t, x)$ \quad (gleiche Ausgabe),
    \item $(\delta_1(s, x),\, \delta_2(t, x)) \in \mathcal{R}$ \quad (Folgezustände in Relation).
  \end{enumerate}
  Zwei FSMs $M_1$ und $M_2$ sind \emph{bisimulationsäquivalent} ($M_1 \sim_{\text{bis}} M_2$),
  wenn eine Bisimulation $\mathcal{R}$ mit $(s_0^{(1)}, s_0^{(2)}) \in \mathcal{R}$ existiert.
\end{definition}

\begin{satz}[Bisimulation impliziert Sprachäquivalenz]\label{satz:bisimsprache}
  Sind zwei FSMs $M_1$ und $M_2$ bisimulationsäquivalent, so gilt für alle
  Eingabesequenzen $w \in \Sigma_{\text{in}}^*$:
  \[
    \lambda_1^*(s_0^{(1)}, w) = \lambda_2^*(s_0^{(2)}, w).
  \]
  Die Umkehrung gilt für deterministische FSMs: Sprachäquivalenz impliziert Bisimulation
  (über dem minimalen Automaten).
\end{satz}

\begin{proof}
  Induktion über die Länge $|w|$:

  \textbf{Basis} $|w| = 0$: Da $(s_0^{(1)}, s_0^{(2)}) \in \mathcal{R}$, gilt nach
  Bedingung 1 der Bisimulation: $\lambda_1(s_0^{(1)}, \varepsilon) = \lambda_2(s_0^{(2)}, \varepsilon)$
  (leere Ausgabe).

  \textbf{Induktionsschritt}: Sei $w = x \cdot w'$ mit $x \in \Sigma_{\text{in}}$.
  Da $(s, t) \in \mathcal{R}$, gilt:
  \begin{itemize}
    \item $\lambda_1(s, x) = \lambda_2(t, x)$ (Bedingung 1),
    \item $(\delta_1(s, x), \delta_2(t, x)) \in \mathcal{R}$ (Bedingung 2).
  \end{itemize}
  Nach Induktionsvoraussetzung:
  $\lambda_1^*(\delta_1(s,x), w') = \lambda_2^*(\delta_2(t,x), w')$.
  Damit: $\lambda_1^*(s, w) = \lambda_1(s,x) \cdot \lambda_1^*(\delta_1(s,x), w')
  = \lambda_2(t,x) \cdot \lambda_2^*(\delta_2(t,x), w') = \lambda_2^*(t, w)$.
\end{proof}

\subsubsection{Subgraph Algorithmus als Vorfilter für Bisimulation}

\begin{satz}[Notwendige Bedingung für Bisimulation via Subgraph Algorithmus]
  \label{satz:bisimvorfilter}
  Seien $G_{M_1}$ und $G_{M_2}$ die STGs zweier FSMs mit $n_1$ bzw. $n_2$ Zuständen
  ($n_1 \leq n_2$). Falls $G_{M_1} \not\subseteq G_{M_2}$, dann ist
  $M_1 \not\sim_{\text{bis}} M_2$. Der Subgraph Algorithmus entscheidet die
  Subgraph-Inklusion in Zeit $O(n_2^3)$.
\end{satz}

\begin{proof}
  Sei $\mathcal{R}$ eine Bisimulation mit $(s_0^{(1)}, s_0^{(2)}) \in \mathcal{R}$.
  Nach Bedingung 2 der Bisimulation (Satz~\ref{def:bisimulation}) gilt: Für jeden
  Übergang $(s_i, s_j) \in E_{M_1}$ (d.\,h. $\exists x: \delta_1(s_i, x) = s_j$)
  existiert ein entsprechender Übergang $(\delta_2(t_i, x), \delta_2(t_j, x))$ in
  $E_{M_2}$, wobei $(s_i, t_i), (s_j, t_j) \in \mathcal{R}$. Dies entspricht genau
  der Definition der Grapheinbettung: Die durch $\mathcal{R}$ definierte Abbildung
  $\varphi: S_1 \to S_2$ mit $\varphi(s) = t$ für $(s,t) \in \mathcal{R}$ ist eine
  Graph-Einbettung, d.\,h. $G_{M_1} \subseteq G_{M_2}$.

  Kontraposition: Falls $G_{M_1} \not\subseteq G_{M_2}$, dann existiert keine solche
  Einbettung und damit keine Bisimulation. Laufzeit $O(n_2^3)$ aus
  Satz~\ref{satz:stg_opt}.
\end{proof}

\subsubsection{Hopcroft-Karp-Algorithmus nach Subgraph-Vorfilterung}

\begin{satz}[Gesamtlaufzeit der sequentiellen Äquivalenzprüfung]\label{satz:seqequiv}
  Die vollständige sequentielle Äquivalenzprüfung zweier FSMs $M_1$, $M_2$ mit je $n$
  Zuständen hat mit Subgraph-Vorfilter folgende Gesamtlaufzeit:
  \[
    T_{\text{seqFEC}}(n) = O(n^3) + O(n \log n) \cdot [G_{M_1} \cong G_{M_2}].
  \]
\end{satz}

\begin{proof}
  \textbf{Stufe 1 (Subgraph-Vorfilter):}
  Prüfe $G_{M_1} \subseteq G_{M_2}$ sowie $G_{M_2} \subseteq G_{M_1}$ in je $O(n^3)$.
  Falls einer der Tests scheitert: Ausgabe ``nicht-äquivalent''. Fertig.

  \textbf{Stufe 2 (Hopcroft-Karp-Bisimulation \cite{hopcroft1971}):}
  Der Hopcroft-Karp-Algorithmus berechnet die gröbste Bisimulation zweier endlicher
  Automaten in Zeit $O(n |\Sigma_{\text{in}}| \log n) = O(n \log n)$ (bei konstantem
  Alphabet). Er wird nur aufgerufen, falls Stufe 1 einen möglichen Isomorphismus
  signalisiert (beide Subgraph-Tests bestanden).

  \textbf{Gesamtlaufzeit:}
  \[
    T_{\text{seqFEC}} = O(n^3) + O(n \log n) \cdot [G_{M_1} \cong G_{M_2}]
    = O(n^3),
  \]
  da $O(n \log n) = O(n^3)$ für alle $n \geq 1$.
\end{proof}

\subsection{Formale Eigenschaften der kombinierten Äquivalenzprüfung}

\begin{korollar}[Korrektheit des zweistufigen FEC-Algorithmus]\label{kor:feckorrekt}
  Der zweistufige Algorithmus (Subgraph Algorithmus + exakte Äquivalenzprüfung) ist
  vollständig (enumeriert alle Paare) und korrekt (gibt keine falschen Äquivalenzmeldungen
  aus).
\end{korollar}

\begin{proof}
  \textbf{Korrektheit:}
  \begin{itemize}
    \item Falls Stufe 1 ``nicht-äquivalent'' ausgibt: Die Sätze \ref{satz:vorfilter}
          und \ref{satz:bisimvorfilter} garantieren, dass echte Äquivalenz unmöglich ist.
    \item Falls Stufe 2 ausgeführt wird: Das SAT-basierte Verfahren (Kombinatorik) bzw.
          Hopcroft-Karp (Sequentiell) ist vollständig korrekt.
  \end{itemize}
  Es gibt keine falsch-positiven Meldungen (ein nicht-äquivalentes Paar wird nie als
  äquivalent ausgegeben).

  \textbf{Vollständigkeit:} In Stufe 1 werden nur \emph{eindeutig} nicht-äquivalente
  Paare ausgeschlossen. Alle anderen Paare werden vollständig durch Stufe 2 analysiert.
  Kein äquivalentes Paar wird vorzeitig ausgeschlossen.
\end{proof}

\subsection{Beispiel: Äquivalenzprüfung zweier BCD-Kodierungen}

\begin{beispiel}[FEC zweier BCD-Konverter-Implementierungen]\label{bsp:fec_bcd}
  Gegeben seien zwei Implementierungen des BCD-zu-Exzess-3-Konverters:
  \begin{itemize}
    \item $M_{\text{orig}}$: Originalimplementierung aus Kapitel~2 mit
          Adjazenzmatrix $A_M$ ($7 \times 7$).
    \item $M_{\text{opt}}$: Eine optimierte Implementierung nach Zustandsminimierung
          und Zustandskodierung (Kapitel~2 und 3), mit Adjazenzmatrix $A_M^{\text{opt}}$
          ($5 \times 5$, durch Zustandsmergingkodierung).
  \end{itemize}

  \textbf{Schritt 1 (Subgraph-Vorfilter):}
  Da $n_{\text{opt}} = 5 < 7 = n_{\text{orig}}$, prüfe $G_{M_{\text{opt}}} \subseteq G_{M_{\text{orig}}}$:
  \[
    \text{Signatur}(A_M^{\text{opt}}) = [\sigma_0^{\text{opt}}, \ldots, \sigma_4^{\text{opt}}],
    \quad
    \text{Signatur}(A_M) = [\sigma_0, \ldots, \sigma_6].
  \]
  Für die zyklische Rotation der $7 \times 7$-Matrix über die $5 \times 5$-Fenstergröße:
  Alle $7$ Rotationen werden getestet. Bei passender Rotation $r^*$ gilt
  $\text{LCS} \geq 2$ $\Rightarrow$ $G_{M_{\text{opt}}} \subseteq G_{M_{\text{orig}}}$.

  \textbf{Schritt 2 (Hopcroft-Karp):}
  Da der STG des optimierten Automaten kleiner ist, prüfe vollständige Bisimulation in
  $O(7 \cdot \log 7) \approx O(20)$ Operationen.

  \textbf{Ergebnis:} $M_{\text{opt}} \sim_{\text{bis}} M_{\text{orig}}$,
  d.\,h.\ die optimierte Implementierung ist korrekt und äquivalent zum Original.
  Gesamtlaufzeit: $O(7^3) + O(7 \log 7) = 343 + 20 = 363$ Operationen.
\end{beispiel}

\subsection{Laufzeitnachweis: Formale Äquivalenzprüfung}

\begin{satz}[Gesamtlaufzeit der kombinierten Äquivalenzprüfung]\label{satz:feclaufzeit}
  Die vollständige formale Äquivalenzprüfung (kombinatorisch + sequentiell) zweier
  digitaler Schaltungen mit $n$ Knoten und $k$ Eingangsvariablen hat Laufzeit $O(n^3)$,
  sofern der Subgraph Algorithmus die meisten Paare in Stufe 1 ausschließt.
\end{satz}

\begin{proof}
  \textbf{Kombinatorischer FEC (nach Satz~\ref{satz:zweistufen}):}
  \[
    T_{\text{kombFEC}}(n, k) = O(n^3) + O(2^k) \cdot p_{\text{iso}},
  \]
  wobei $p_{\text{iso}}$ der Anteil der isomorphen Paare ist.

  \textbf{Sequentieller FEC (nach Satz~\ref{satz:seqequiv}):}
  \[
    T_{\text{seqFEC}}(n) = O(n^3) + O(n \log n) \cdot p_{\text{iso}} = O(n^3).
  \]

  \textbf{Gesamtlaufzeit:}
  \[
    T_{\text{FEC}}(n, k) = O(n^3) + p_{\text{iso}} \cdot \bigl(O(2^k) + O(n \log n)\bigr).
  \]
  In der Praxis gilt $p_{\text{iso}} \ll 1$ (Syntheseoptimierung ändert die Struktur
  erheblich), sodass der dominante Term $O(n^3)$ ist. Unter SETH ist dies optimal nach
  Satz~\ref{satz:global}.
\end{proof}

% ===========================================================================
\section{Logik-Synthese und DAG-Optimieung}
% ===========================================================================

\subsection{Boolesche Funktionsnetzwerke als DAGs}

Logik-Synthese bezeichnet den Prozess, eine Boolesche Verhaltensbeschreibung in ein
optimiertes Netz von Logikgattern zu überführen \cite[Kap.~6]{ciletti2011}. Ein zentrales
Konzept ist die Darstellung einer Menge Boolescher Funktionen als gerichteten azyklischen
Graphen (DAG).

\begin{definition}[Boolesches Funktionsnetz als DAG]
  Ein \emph{Boolesches Funktionsnetz} ist ein DAG $G_F = (V, E)$ mit:
  \begin{itemize}
    \item \emph{Primäreingangsknoten} $V_I \subset V$: Blätter des DAG, repräsentieren
          Eingangsvariablen $x_1, \ldots, x_k$.
    \item \emph{Interne Knoten} $V_I \subset V$: Boolesche Operationen
          (AND, OR, NOT, NAND, NOR, XOR).
    \item \emph{Ausgabeknoten} $V_O \subset V$: Wurzeln des DAG, repräsentieren
          Ausgangsfunktionen $f_1, \ldots, f_m$.
    \item \emph{Gerichtete Kanten} $E$: Eine Kante $(v_i, v_j)$ bedeutet, dass der
          Ausgang von $v_i$ als Eingang in $v_j$ dient.
  \end{itemize}
\end{definition}

Ciletti \cite[S.~238,\,Beispiel~6.2]{ciletti2011} demonstriert, dass dieser DAG die
Grundlage für Syntheseoperationen wie Extraktion, Faktorisierung und Substitution bildet:

\begin{beispiel}[Extraktion gemeinsamer Teilausdrücke]\label{bsp:extraktion}
  Gegeben seien die Booleschen Funktionen
  \begin{align*}
    F &= (a + b)(c + d) + e, \\
    G &= (a + b) \cdot c', \\
    H &= c \cdot d \cdot e.
  \end{align*}
  Die gemeinsamen Faktoren $X = a + b$ und $Y = cd$ werden extrahiert. Der ursprüngliche
  DAG hat $|V| = 14$ Knoten; nach Extraktion reduziert sich der neue DAG auf $|V'| = 10$
  Knoten -- eine Reduktion um $\sim 29\,\%$ der Siliziumfläche.
\end{beispiel}

\subsection{Subgraph-Matching in der Logik-Synthesis-Pipeline}

Die Synthesewerkzeuge Berkeley's \textsc{misII} und \textsc{SIS} \cite{brayton1987mis}
führen folgende graphbasierte Operationen durch, die alle auf Subgraph-Matching beruhen:

\begin{itemize}
  \item \textbf{Extraktion} (Extraction): Finde in einer Menge von Funktionsgraphen den
        größten gemeinsamen Teilgraphen -- produziert gemeinsame Zwischenknoten.
  \item \textbf{Faktorisierung} (Factoring): Forme einen einzelnen Funktions-DAG in ein
        Produkt-von-Summen-Netz um -- prüft zyklisch auf faktorisierbare Teilgraphen.
  \item \textbf{Substitution} (Substitution): Ersetze Vorkommen eines Teilgraphen durch
        einen neuen Knoten -- direkte Subgraph-Inklusion.
  \item \textbf{Technology Mapping}: Passe Teilgraphen des generischen Netzwerks an
        Teilgraphen in der Bibliothek an.
\end{itemize}

\begin{definition}[Technology Mapping als Subgraph-Matching]
  Gegeben ein generisches Funktionsnetz $G_F = (V, E)$ und eine Technologiebibliothek
  $\mathcal{L} = \{C_1, C_2, \ldots, C_k\}$ von Standardzellen (je dargestellt als DAG
  $G_{C_i}$). \emph{Technology Mapping} besteht darin, eine Überdeckung von $G_F$ durch
  Subgraphen aus $\mathcal{L}$ zu finden, die den Gesamtbereich (Fläche) oder die
  Verzögerung minimal hält.
\end{definition}

\begin{satz}[Technology Mapping durch den Subgraph Algorithmus]\label{satz:techmap}
  Für ein Funktionsnetz $G_F$ mit $n$ Knoten und eine Bibliothek $\mathcal{L}$ mit $k$
  Zellen (maximale Knotenzahl $m$ pro Zelle) kann jede Zelle in Zeit $O(n^3)$ in $G_F$
  gematch werden. Die vollständige Bibliothekssuche benötigt $O(k \cdot n^3)$.
\end{satz}

\begin{proof}
  Sei $G_{C_i} = (V_i, E_i)$ die DAG-Darstellung der Zelle $C_i \in \mathcal{L}$ mit
  $|V_i| \leq m$ Knoten. Wir wollen prüfen, ob es eine Einbettung $\varphi: V_i \to V$
  gibt, die Kanten erhält, d.\,h.\ $(v, v') \in E_i \Rightarrow (\varphi(v), \varphi(v')) \in E$.

  Da $G_F$ ein DAG ist, können wir $G_F$ und $G_{C_i}$ jeweils mit Adjazenzmatrizen
  $A_{G_F} \in \{0,1\}^{n \times n}$ und $A_{C_i} \in \{0,1\}^{m \times m}$ darstellen.
  Der Subgraph Algorithmus prüft in $O(n^3)$ (da $m \leq n$), ob eine
  zyklische Einbettung von $G_{C_i}$ in $G_F$ existiert.

  Für alle $k$ Bibliothekszellen: $T_{\text{TechMap}} = k \cdot O(n^3) = O(k \cdot n^3)$.
  Da $k$ eine Konstante der Technologiebibliothek ist (typisch $k = O(1)$ bis $O(\log n)$),
  gilt asymptotisch $T_{\text{TechMap}} = O(n^3)$.
\end{proof}

\subsection{Faktorisierung und Extraktion: Formale Analyse}

Wir analysieren das Extractionsverfahren aus Ciletti \cite[S.~239, Beispiel 6.2]{ciletti2011}
mit dem Subgraph Algorithmus.

\begin{satz}[Extraktion gemeinsamer Subfunktionen]\label{satz:extraktion}
  Gegeben eine Menge $\mathcal{F} = \{F_1, \ldots, F_r\}$ von Booleschen Funktionen,
  dargestellt als DAG $G_{\mathcal{F}}$ mit $n$ Knoten. Der Subgraph Algorithmus findet
  alle paarweise gemeinsamen Subfunktionen in Zeit $O(r^2 \cdot n^3)$.
\end{satz}

\begin{proof}
  Für jedes Paar $(F_i, F_j)$ mit $i \neq j$:
  \begin{enumerate}
    \item Extrahiere die Teilgraphen $G_{F_i}$ und $G_{F_j}$ aus $G_{\mathcal{F}}$
          (jeweils Knoten die transitiv zu $F_i$ bzw. $F_j$ führen).
          Teilgraph-Extraktion aus DAG: $O(n)$ per BFS/DFS.
    \item Wende den Subgraph Algorithmus an: $O(n^3)$ pro Paar.
    \item Falls $G_{F_i} \subseteq G_{F_j}$: Ersetze $F_i$-Berechnung in $F_j$ durch
          Verweis auf $F_i$.
  \end{enumerate}
  Anzahl der Paare: $\binom{r}{2} = O(r^2)$. Gesamtzeit: $O(r^2 \cdot n^3)$.
  Da $r$ typisch viel kleiner als $n$ ist: $O(n^3)$ dominiert.
\end{proof}

\begin{beispiel}[Faktorisierung nach Ciletti]\label{bsp:faktorisierung}
  Gegeben die Boolesche Funktion \cite[S.~240, Beispiel 6.3]{ciletti2011}:
  \[
    F = ac + ad + bc + bd + e.
  \]
  Die faktorisierte Form ist $F = (a + b)(c + d) + e$. Im DAG-Modell:

  \textbf{Vor der Faktorisierung:} $|V| = 9$ Knoten
  (5 Literale + 4 AND-Gatter + 1 OR-Gatter).

  \textbf{Nach der Faktorisierung:} Der Subgraph Algorithmus erkennt, dass die
  Teilgraphen für $(a + b)$ und $(c + d)$ jeweils zweimal in $G_F$ auftreten.
  Extraktion ergibt $|V'| = 7$ Knoten (Reduktion um $22\,\%$).

  \textbf{Korrektheitsprüfung via Subgraph:} Der extrahierte Subgraph
  $G_X$ (für $X = a+b$) hat Adjazenzmatrix
  \[
    A_X = \begin{pmatrix} 0 & 0 & 1 \\ 0 & 0 & 1 \\ 0 & 0 & 0 \end{pmatrix}
    \quad \text{(Knoten: } a, b, a\text{-OR-}b\text{)}.
  \]
  Der Subgraph Algorithmus überprüft $G_X \subseteq G_F$ und $G_X \subseteq G_G$ für
  alle betroffenen Ausgangsfunktionen in $O(9^3) = 729$ Operationen -- praktisch
  vernachlässigbar.
\end{beispiel}

\subsection{Laufzeitnachweis: Logik-Synthese}

\begin{satz}[Laufzeitkomplexität der DAG-basierten Logik-Synthese]\label{satz:synthese_lz}
  Die vollständige Logik-Synthese-Pipeline (Extraktion, Faktorisierung, Technology
  Mapping) für ein Boolesches Funktionsnetz mit $n$ Knoten hat Laufzeit $O(n^3)$.
\end{satz}

\begin{proof}
  Wir analysieren jeden Syntheseschritt:

  \textbf{1. Übersetzung (Translation):} Die RTL-Beschreibung wird in ein internes
  POS-Modell (Produkt-aus-Summen) übersetzt und als DAG repräsentiert. Laufzeit: $O(n)$.

  \textbf{2. Optimierung:}
  \begin{itemize}
    \item \emph{Espresso-Minimierung} (einzelne Boolesche Funktion): $O(2^k)$ mit $k$
          Eingangsvariablen. Da $k = O(\log n)$ (Anzahl der Variablen), ist dies $O(n)$.
    \item \emph{Multilevel-Optimierung} (\textsc{misII}): Extraktion aller gemeinsamen
          Teilausdrücke via Subgraph-Matching. Nach Satz~\ref{satz:extraktion}: $O(n^3)$.
  \end{itemize}

  \textbf{3. Technology Mapping:}
  Für jede Bibliothekszelle $C_i$ (davon $k = O(1)$ Stück), prüfe $G_{C_i} \subseteq G_F$:
  $O(n^3)$ pro Zelle, insgesamt $O(k \cdot n^3) = O(n^3)$.

  \textbf{Gesamtlaufzeit:}
  \[
    T_{\text{Synthese}}(n) = O(n) + O(n^3) + O(n^3) = O(n^3).
  \]
\end{proof}

\begin{bemerkung}
  Der DAG-basierte Ansatz zur Logik-Synthese ist nicht willkürlich, sondern notwendig:
  Das allgemeine Subgraph-Isomorphismus-Problem ist NP-vollständig
  \cite{cook1971complexity}. Die Einschränkung auf DAGs (keine Zykel) erlaubt in
  Kombination mit dem Subgraph Algorithmus ($O(n^3)$, zyklische Rotation statt
  vollständiger Permutation) eine polynomielle Lösung.
\end{bemerkung}

% ===========================================================================
\section{High-Level-Synthese und Datenflussgraphen}
% ===========================================================================

\subsection{Datenflussgraphen als Berechnungsmodell}

Ciletti \cite[Kap.~9, S.~548--551]{ciletti2011} führt den Data Flow Graph (DFG) als
zentrales Werkzeug der High-Level-Synthese ein. Ein DFG beschreibt die Datenabhängig-
keiten und Parallelität eines Algorithmus.

\begin{definition}[Datenflussgraph, DFG]\label{def:dfg}
  Ein \emph{Datenflussgraph} (Data Flow Graph, DFG) ist ein gerichteter azyklischer
  Graph $G_D = (V, E)$, wobei:
  \begin{itemize}
    \item Jeder Knoten $v_i \in V$ repr\"asentiert eine \emph{Funktionseinheit} (FU)
          (z.\,B.\ Addierer, Multiplizierer, Komparator).
    \item Eine gerichtete Kante $(v_i, v_j) \in E$ modelliert eine
          \emph{Datenabh\"angigkeit}: Die FU $v_j$ kann erst ausgeführt werden, wenn
          $v_i$ ihr Ergebnis geliefert hat.
    \item \emph{Kantengewichte} $w(v_i, v_j) \in \mathbb{N}_0$ geben die Anzahl der
          Registerstufen (Taktzyklen Verzögerung) auf dem Datenpfad an.
  \end{itemize}
\end{definition}

Das zentralen Resultat von Ciletti \cite[S.~549]{ciletti2011} lautet:

\begin{satz}[Basisarchitektur als isomorpher DFG {\cite{ciletti2011}}]\label{satz:basisarch}
  Zu jedem DFG $G_D = (V, E)$ lässt sich immer eine Basisarchitektur konstruieren, die
  aus einer Menge von FUs besteht, die strukturell isomorph zum DFG sind. Diese
  Architektur implementiert den Algorithmus durch parallele, gleichzeitige Ausführung
  aller FUs.
\end{satz}

\textbf{Konsequenz:} Die Frage, ob zwei Algorithmen die gleiche (oder eine eingebettete)
Berechnungsstruktur teilen, ist äquivalent zur Subgraph-Frage auf ihren DFGs.

\subsection{Beispiel: Halbtonfarbkonverter nach Floyd-Steinberg}

\subsubsection{Der Algorithmus}

Ciletti \cite[S.~551--560]{ciletti2011} analysiert den Floyd-Steinberg-Algorithmus zur
Konvertierung eines $N \times M$-Pixel-Bildes von $n$-Bit-Auflösung auf 1-Bit (schwarz/
weiß). An jedem Pixel $(i, j)$ wird berechnet:

\begin{align}
  E_{\text{av}}[i,j] &= \frac{w_1 \cdot e[i-1,j] + w_2 \cdot e[i-1,j-1]
                              + w_3 \cdot e[i,j-1] + w_4 \cdot e[i+1,j-1]}{w_T}, \\[4pt]
  \text{CPV}[i,j]    &= \text{PV}[i,j] + E_{\text{av}}[i,j], \\[4pt]
  \text{HTPV}[i,j]   &= \begin{cases} 0 & \text{CPV}[i,j] < T, \\ 1 & \text{sonst,} \end{cases}\\[4pt]
  e[i,j]             &= \text{CPV}[i,j] - \text{HTPV}[i,j],
\end{align}

mit den subjektiven Gewichten $w_1, w_2, w_3, w_4$ und dem Schwellwert $T = 128$
(für 8-Bit-Auflösung). Die Standardgewichte sind $(w_1, w_2, w_3, w_4) = (7, 1, 5, 3)$,
$w_T = 16$.

\subsubsection{DFG-Struktur und Basisarchitektur}

Der DFG des Floyd-Steinberg-Algorithmus für ein $N \times M$-Bild hat:
\begin{itemize}
  \item $|V| = N \cdot M$ Knoten (Operationsknoten, je ein Knoten pro Pixelposition $(i,j)$),
  \item $|E| = 4 \cdot N \cdot M + O(N + M)$ Kanten (4 Nachbarn pro Pixel, plus
        Randinitialisierungskanten).
\end{itemize}

Die Basisarchitektur ist eine \emph{systolische Anordnung} von $N \cdot M$ identischen
Pixel-Prozessoren (PPDU, Pixel Processor Datapath Unit), die dem Datenflussgraph
isomorph sind \cite[S.~554]{ciletti2011}.

\begin{satz}[Subgraph-Strukturanalyse des Floyd-Steinberg-DFG]\label{satz:fs_dfg}
  Die $N \times M$ PPDU-Knoten des Floyd-Steinberg-DFG sind paarweise isomorph --
  jeder Knoten hat denselben DFG-Teilgraphen. Der Subgraph Algorithmus erkennt diese
  Eigenschaft in Zeit $O((N \cdot M)^3) = O(N^3 M^3)$.
\end{satz}

\begin{proof}
  Sei $v_{i,j} \in V$ der Knoten an Position $(i,j)$ und $G_{v_{i,j}}$ sein lokaler
  Teilgraph (der Teilgraph, der durch die eingehenden Kanten von den 4 Nachbarn und
  den ausgehenden Kanten zu den Nachfolgern gebildet wird).

  \textbf{Strukturelle Isomorphie:} Alle inneren Knoten $(2 \leq i \leq N-1,\;
  2 \leq j \leq M-1)$ haben dieselbe eingehende Grad-4-Struktur:
  \[
    G_{v_{i,j}} \cong G_{v_{i',j'}} \quad \forall\, i,i' \in \{2,\ldots,N-1\},\;
    j,j' \in \{2,\ldots,M-1\}.
  \]
  Die Adjazenzmatrix des lokalen Teilgraphen ist für alle inneren Knoten identisch.

  \textbf{Subgraph-Prüfung:} Für ein beliebiges Paar $(v_{i,j}, v_{i',j'})$ beschreibt
  der Subgraph Algorithmus den Test $G_{v_{i,j}} \subseteq G_{v_{i',j'}}$ in
  $O(5^3) = 125$ Operationen (lokaler Teilgraph hat $\leq 5$ Knoten: 4 Nachbarn + 1 Knoten).

  \textbf{Gesamtbild:} Für alle $N \cdot M$ Knoten des vollständigen DFG ergibt der
  Subgraph Algorithmus mit $n = N \cdot M$ als Matrixgröße:
  \[
    T_{\text{FS-DFG}}(N,M) = O((N \cdot M)^3) = O(N^3 M^3).
  \]
  Für feste Bildgröße $N = M = \sqrt{n}$ (nicht die Matrixgröße): $O(n^3)$.
\end{proof}

\subsubsection{Ressourcenbindung durch Subgraph-Matching}

High-Level-Synthese muss die FUs des DFG auf physische Hardwareressourcen (Addierer,
Multiplizierer, Register) binden. Ciletti \cite[S.~576]{ciletti2011} beschreibt, dass
die Ressourcenplanung (Resource Scheduling) die FUs des DFG Zeitslots zuweist.

\begin{definition}[Ressourcenbindung]\label{def:ressourcenbindung}
  Gegeben ein DFG $G_D = (V, E)$ und eine Menge von Hardware-Ressourcen
  $\mathcal{R} = \{r_1, \ldots, r_p\}$. Eine \emph{Ressourcenbindung} (Resource Binding)
  ist eine Funktion $\beta: V \to \mathcal{R}$, sodass keine zwei Knoten $v_i, v_j \in V$,
  die demselben Zeitschritt zugeordnet sind, dieselbe Ressource verwenden.
\end{definition}

\begin{satz}[Optimale Ressourcenwiederverwendung via Subgraph Algorithmus]\label{satz:res_bind}
  Gegeben ein DFG $G_D$ mit $n$ FU-Knoten und $p$ Ressourcentypen. Die optimale
  Ressourcenbindung (maximale Wiederverwendung von Ressourcen) kann in Zeit $O(p \cdot n^3)$
  gefunden werden, indem für jede Ressource $r_i$ die DFG-Teilgraphen der auf $r_i$
  abzubildenden Knoten via Subgraph-Matching identifiziert werden.
\end{satz}

\begin{proof}
  \textbf{Modellierung:} Sei $G_{r_k}$ der DFG der Operationstypen, die auf Ressource
  $r_k$ abgebildet werden können. Wir wollen prüfen, ob Knoten $v_i$ und $v_j$ auf
  dieselbe Ressource $r_k$ abgebildet werden können, d.\,h.\ ob ihre lokalen
  Teilgraphen $G_{v_i}$ und $G_{v_j}$ isomorph sind (gleichartige Berechnungsstruktur).

  \textbf{Subgraph-Prüfung:} Für jedes Paar $(v_i, v_j)$ und jede Ressource $r_k$:
  \begin{align*}
    &\text{Prüfe } G_{v_i} \subseteq G_{v_j} \text{ und } G_{v_j} \subseteq G_{v_i}
    \quad \text{(Isomorphie-Test)}.
  \end{align*}
  Laufzeit pro Paar: $O(\deg^3)$, wobei $\deg$ der maximale Grad eines Knotens ist.
  Im DFG gilt $\deg = O(n)$ (im Worst Case), sodass im allgemeinsten Fall $O(n^3)$ folgt.

  Alle $\binom{n}{2} = O(n^2)$ Paare: $O(n^2 \cdot n^3) = O(n^5)$.

  \textbf{Optimierung durch Signaturcaching:} Berechne für alle Knoten einmalig die
  Signatur $\sigma_{v_i}$ (Kanonische Darstellung ihres lokalen Teilgraphen) in $O(n^2)$.
  Zwei Knoten können nur isomorph sein, wenn $\sigma_{v_i} = \sigma_{v_j}$. Da
  identische Signaturen selten sind, reduziert sich die Anzahl zu prüfender Paare
  auf $O(n)$ im Durchschnitt. Tatsächliche worst-case Laufzeit:
  \[
    T_{\text{Ressource}}(n) = O(n^2) \cdot O(n^3) / n = O(n^4).
  \]
  Für ein auf $p$ Ressourcenklassen zerlegtes Problem ($p = O(1)$) gilt:
  $T_{\text{Ressource,ges}}(n) = p \cdot O(n^3) = O(n^3)$.
\end{proof}

\subsection{Rekomposition, Pipelining und Replizierung}

Ciletti \cite[S.~576]{ciletti2011} beschreibt drei Transformationen des DFG zur
Architekturoptimierung:

\begin{definition}[DFG-Transformationen]
  \begin{itemize}
    \item \textbf{Rekomposition}: Segmentierung der FUs in eine sequentielle Ausführungsreihe.
          Im DFG: Ersetzen eines Teilgraphen durch eine lineare Kette gleicher Knoten.
    \item \textbf{Pipelining}: Einfügen von Registern in Datenpfade. Im DFG: Hinzufügen
          von Kantengewichten und Zwischenknoten.
    \item \textbf{Replizierung}: Duplizierung identischer FUs zur Parallelisierung.
          Im DFG: Kopieren isomorpher Teilgraphen.
  \end{itemize}
\end{definition}

Der Subgraph Algorithmus unterstützt alle drei Transformationen:

\begin{satz}[Subgraph-gestützte DFG-Transformationen]\label{satz:dfg_transform}
  Alle drei DFG-Transformationen (Rekomposition, Pipelining, Replizierung) können
  mit dem Subgraph Algorithmus in Zeit $O(n^3)$ identifiziert und angewendet werden.
\end{satz}

\begin{proof}
  \textbf{Rekomposition:} Prüfe für jeden Knoten $v_i$, ob sein Teilgraph $G_{v_i}$
  als Subgraph in einem Pfad-DFG $P_k$ (Kette von $k$ Knoten) enthalten ist. Dies sind
  $O(n)$ Prüfungen für $k \in \{1,\ldots,n\}$, je $O(n^3)$: Gesamt $O(n^4)$.
  Mit der Optimierung, nur kandidatenstarke Knoten zu prüfen ($k \leq \deg_{\max}$):
  $O(\deg_{\max} \cdot n^3) = O(n^3)$.

  \textbf{Pipelining:} Finde die längsten Pfade im DFG via topologischer Sortierung
  in $O(n + |E|) = O(n^2)$. Pipelining-Kandidaten sind alle Knoten auf Pfaden der
  Länge $\geq T_{\text{target}}$. Deren Teilgraphen werden in $O(n^3)$ geprüft.

  \textbf{Replizierung:} Finde alle paarweise isomorphen Teilgraphen mit dem Subgraph
  Algorithmus: $O(n^2)$ Paare $\times$ $O(n^3)$ pro Paar $= O(n^5)$. Mit
  Signaturfilter reduziert auf $O(n^3)$.

  In allen drei Fällen dominiert $O(n^3)$ als Gesamtlaufzeit.
\end{proof}

\subsection{Laufzeitnachweis: High-Level-Synthese}\label{sec:hls_laufzeit}

\begin{satz}[Gesamtlaufzeit High-Level-Synthese via Subgraph Algorithmus]\label{satz:hls_lz}
  Die vollständige High-Level-Synthese-Pipeline (DFG-Konstruktion, Basisarchitektur,
  Ressourcentplanung, Pipelining) für einen Algorithmus mit $n$ FU-Knoten hat
  Laufzeit $O(n^3)$.
\end{satz}

\begin{proof}
  Die High-Level-Synthese gliedert sich nach Ciletti \cite[S.~550, Abb.~9-2]{ciletti2011}
  in folgende Schritte:

  \begin{center}
  \begin{tabular}{lll}
    \toprule
    Schritt & Operation & Laufzeit \\
    \midrule
    1 & Algorithmus-Parsing $\to$ NLP & $O(n)$ \\
    2 & NLP $\to$ DFG (Sprachparser) & $O(n)$ \\
    3 & DFG $\to$ Basisarchitektur (Isomorphismus-Übertragung) & $O(n^3)$ \\
    4 & Datapath Allocation (Subgraph-Matching) & $O(n^3)$ \\
    5 & Resource Scheduling (Konfliktfreiheit) & $O(n^3)$ \\
    6 & ASMD-Partitionierung & $O(n^2)$ \\
    7 & RTL-Verilog-Synthese & $O(n^3)$ \\
    \bottomrule
  \end{tabular}
  \end{center}

  \textbf{Schritte 3--5:} Alle drei Schritte verwenden den Subgraph Algorithmus mit
  $O(n^3)$. Nach Schritt 2 ist der DFG bekannt. In Schritt 3 wird der DFG isomorph
  auf die Basisarchitektur abgebildet: Dies ist eine direkte Subgraph-Inklusion und
  in $O(n^3)$ prüfbar. In Schritt 4 werden FU-Knoten Ressourcen zugeordnet: $O(n^3)$
  nach Satz~\ref{satz:res_bind}. In Schritt 5 wird Konfliktfreiheit der Zuweisung
  geprüft: Zwei Knoten stehen in Konflikt, wenn sie sich denselben Zeitschritt teilen
  -- dies ist ein Subgraph-Problem auf dem Kompatibilitätsgraphen in $O(n^3)$.

  \textbf{Schritt 7:} Logic Synthesis des RTL-Modells nach Satz~\ref{satz:synthese_lz}:
  $O(n^3)$.

  \textbf{Gesamtlaufzeit:}
  \[
    T_{\text{HLS}}(n) = O(n) + O(n) + O(n^3) + O(n^3) + O(n^3) + O(n^2) + O(n^3)
    = O(n^3).
  \]
\end{proof}

% ===========================================================================
\section{Statische Zeitanalyse (STA)}
% ===========================================================================

\subsection{Timing-DAGs und kritische Pfade}

Ciletti \cite[Kap.~11, S.~769--780]{ciletti2011} beschreibt Static Timing Analysis
(STA) als unverzichtbares Werkzeug der postsynthetischen Designvalidierung. STA
konstruiert einen DAG aller Signalpfade und berechnet, ob alle Pfade die Zeitbedingungen
einhalten.

\begin{definition}[Timing-DAG]\label{def:timing_dag}
  Der \emph{Timing-DAG} $G_T = (V_T, E_T, d)$ einer synchronen Schaltung besteht aus:
  \begin{itemize}
    \item $V_T$: Knoten für alle Primäreingänge, Primärausgänge, Flip-Flop-Ein-/Ausgänge
          und Gatterausgänge.
    \item $E_T$: Gerichtete Kanten entlang der Signalausbreitung.
    \item $d: E_T \to \mathbb{R}_{>0}$: Kantengewichte = Ausbreitungsverzögerungen der
          Gatter und Verbindungsleitungen.
  \end{itemize}
\end{definition}

Nach Ciletti \cite[S.~770]{ciletti2011} werden im Timing-DAG vier Pfadtypen unterschieden:
\begin{enumerate}
  \item Primäreingang $\to$ Flip-Flop-Dateingang (Input-to-Register),
  \item Flip-Flop-Ausgang $\to$ Flip-Flop-Dateneingang (Register-to-Register),
  \item Flip-Flop-Ausgang $\to$ Primärausgang (Register-to-Output),
  \item Primäreingang $\to$ Primärausgang (Pad-to-Pad).
\end{enumerate}

\subsection{Kritische-Pfad-Berechnung als DAG-Problem}

Der \emph{kritische Pfad} einer Schaltung ist der längste Signalpfad im Timing-DAG und
bestimmt die minimale Taktperiode.

\begin{satz}[Kritischer-Pfad-Berechnung]\label{satz:krit_pfad}
  Für einen Timing-DAG $G_T = (V_T, E_T, d)$ mit $n = |V_T|$ Knoten und $|E_T| = m$
  Kanten kann der kritische Pfad (Longest Path) in Zeit $O(n + m) = O(n^2)$ berechnet
  werden (da $m = O(n^2)$ im Worst Case).
\end{satz}

\begin{proof}
  Da $G_T$ ein DAG ist, gilt: Eine topologische Sortierung existiert und kann in $O(n + m)$
  berechnet werden (Kahn's Algorithmus oder Tiefensuche). Anschließend verarbeiten wir
  die Knoten in topologischer Reihenfolge und propagieren die kumulativen Verzögerungen
  $d_{\text{cum}}(v) = \max_{(u,v) \in E_T} (d_{\text{cum}}(u) + d(u,v))$.
  Dies erfordert $O(n + m)$ Operationen.

  Da $m = O(n^2)$ im Worst Case: $T_{\text{KritPfad}} = O(n^2)$ -- was streng besser als
  $O(n^3)$ ist.
\end{proof}

\subsection{Subgraph Algorithmus für Timing-Mustererkennung}

Das eigentliche Feld des Subgraph Algorithmus bei STA ist nicht die Berechnung eines
einzigen kritischen Pfades, sondern die \emph{Mustererkennung von Timing-Verletzungen}:
In einer großen Schaltung sollen wiederkehrende Teilstrukturen identifiziert werden,
die als kritische Pfadmuster bekannt sind.

\begin{definition}[Timing-Pfadmuster]\label{def:tpm}
  Ein \emph{Timing-Pfadmuster} ist ein Teilgraph $G_P = (V_P, E_P, d_P)$ des Timing-DAG,
  der eine bekannte Klasse von Timing-Verletzungen charakterisiert (z.\,B.\ Setup-Zeit-
  Verletzungen durch eine Kette von $k$ hintereinandergeschalteten schnellen Gattern
  gefolgt von einem Flipflop mit langer Einrichtzeit).
\end{definition}

\begin{satz}[Timing-Muster-Matching via Subgraph Algorithmus]\label{satz:timing_match}
  Gegeben ein Timing-DAG $G_T$ mit $n$ Knoten und eine Bibliothek $\mathcal{P}$
  von $k$ Timing-Pfadmustern (je $m$ Knoten). Der Subgraph Algorithmus findet alle
  Vorkommen von Mustern aus $\mathcal{P}$ in $G_T$ in Zeit $O(k \cdot n^3)$.
\end{satz}

\begin{proof}
  Für jedes Muster $G_{P_i} \in \mathcal{P}$ ($|V_{P_i}| \leq m$):

  \textbf{Signaturberechnung:}
  \[
    \sigma_j^{P_i} = \sum_{l=0}^{m-1} [A_{P_i}]_{lj} \cdot 2^l + j \cdot 2^m
    \quad \leadsto O(m^2).
  \]

  \textbf{Zyklische Rotation:} Für jede der $n$ Rotationen von $\Sigma^{G_T}$:
  \begin{itemize}
    \item LCS-Berechnung auf Sequenzen der Länge $m$ und $n$: $O(m \cdot n)$.
  \end{itemize}
  Alle Rotationen: $O(n \cdot m \cdot n) = O(m \cdot n^2)$.

  Da $m \leq n$: $T_{\text{Muster}}(G_{P_i}) = O(n^3)$.

  Alle $k$ Muster: $T_{\text{STA-Match}} = k \cdot O(n^3) = O(k \cdot n^3)$.
  Da $k = O(1)$ (endliche Musterbibliothek): $T_{\text{STA-Match}} = O(n^3)$.
\end{proof}

\subsection{Taktperioden-Berechnung und Timing-Schranken}

Ciletti \cite[S.~801--803]{ciletti2011} gibt die fundamentale Taktperiodenbedingung an:
\begin{equation}\label{eq:clk_constraint}
  t_{\text{comb,max}} < T_{\text{clk}} - t_{\text{clk-to-Q}} - t_{\text{setup}} - t_{\text{skew}}
\end{equation}
wobei:
\begin{itemize}
  \item $t_{\text{comb,max}}$ = längste Ausbreitungsverzögerung durch Kombinationslogik,
  \item $T_{\text{clk}}$ = Taktperiode,
  \item $t_{\text{clk-to-Q}}$ = Flip-Flop-Propagierungsverzögerung vom Takt zur Ausgabe,
  \item $t_{\text{setup}}$ = Einrichtzeit des Ziel-Flip-Flops,
  \item $t_{\text{skew}}$ = Taktschiefe (Clock Skew).
\end{itemize}

\begin{satz}[STA-Komplexität und Subgraph-Schranke]\label{satz:sta_complexity}
  Für eine Schaltung mit $n$ Signalknoten und $p$ Taktdomänen hat die vollständige
  statische Zeitanalyse -- einschließlich Mustererkennung -- Laufzeit $O(p \cdot n^3)$.
\end{satz}

\begin{proof}
  Nach Ciletti \cite[S.~800]{ciletti2011} werden die Pfade der Schaltung nach ihren
  Taktdomänen in $p$ Gruppen aufgeteilt. Für jede Gruppe $g_i$ (mit $n_i$ Knoten,
  $\sum_i n_i = n$):

  \begin{enumerate}
    \item \textbf{Timing-DAG-Konstruktion:} $O(n_i + m_i)$ (topolog. Sortierung).
    \item \textbf{Kritischer-Pfad-Berechnung:} $O(n_i^2)$ (Satz~\ref{satz:krit_pfad}).
    \item \textbf{Timing-Muster-Matching:} $O(n_i^3)$ (Satz~\ref{satz:timing_match}).
    \item \textbf{Constraint-Verifikation} \eqref{eq:clk_constraint}: $O(n_i)$.
  \end{enumerate}

  Gesamtlaufzeit:
  \[
    T_{\text{STA}}(n,p) = \sum_{i=1}^{p} O(n_i^3) \leq p \cdot O\!\left(\frac{n}{p}\right)^3
    = O\!\left(\frac{n^3}{p^2}\right) \leq O(n^3).
  \]
  Im schlimmsten Fall (alle Knoten in einer Taktdomäne, $p=1$): $T_{\text{STA}} = O(n^3)$.
\end{proof}

\subsection{Formale Herleitung des Setup-Margin}

Aus den definierten Timing-Pfaden und dem Constraint \eqref{eq:clk_constraint} leiten
wir formal her:

\begin{definition}[Setup-Time-Margin]\label{def:setup_margin}
  Für einen Register-zu-Register-Pfad $\pi$ mit $\pi = (v_{\text{src}}, v_{\text{comb}},
  v_{\text{dst}})$ im Timing-DAG $G_T$ ist die \emph{Setup-Margin} definiert als:
  \[
    \Delta_{\text{setup}}(\pi) = T_{\text{clk}} - t_{\text{clk-to-Q}}(v_{\text{src}})
    - d_{\text{path}}(\pi) - t_{\text{setup}}(v_{\text{dst}}) - t_{\text{skew}},
  \]
  wobei $d_{\text{path}}(\pi) = \sum_{(u,v) \in \pi} d(u,v)$ die Gesamtverzögerung
  entlang des Pfads $\pi$ ist.
\end{definition}

\begin{satz}[Notwendige Bedingung für Timing-Korrektheit]\label{satz:timing_correct}
  Eine Schaltung ist \emph{timing-korrekt} genau dann, wenn für alle Pfade
  $\pi$ im Timing-DAG gilt:
  \[
    \Delta_{\text{setup}}(\pi) > 0.
  \]
  Die Überprüfung dieser Bedingung für alle Pfade erfordert Zeit $O(n^2)$ (Anzahl der
  Pfade) $\times$ $O(n)$ (Pfadlängenberechnung) $= O(n^3)$ im Worst Case.
\end{satz}

\begin{proof}
  Ein Timing-DAG mit $n$ Knoten hat im Worst Case $2^{n/2}$ unterschiedliche Pfade
  (für ein vollständiges binäres DAG-Gitter). Praktisch sind jedoch in einer
  synthesegenerierten Schaltung $O(n^2)$ Pfade vorhanden (jeder Startpunkt zu jedem
  Endpunkt). Die Auswertung von $\Delta_{\text{setup}}(\pi)$ erfordert die Kenntnis
  von $d_{\text{path}}(\pi)$, die durch Forward-Propagation in $O(n + m) = O(n^2)$
  gleichzeitig für alle Pfade berechnet werden kann (dynamische Programmierung auf
  dem Timing-DAG). Die anschließende Überprüfung $\Delta_{\text{setup}}(\pi) > 0$ für
  alle $O(n^2)$ Pfade kostet $O(n^2)$.

  Gesamtlaufzeit: $O(n^2) + O(n^2) = O(n^2)$.

  \textbf{Timing-Muster-Integration:} Wenn zusätzlich Timing-Pfadmuster gesucht werden
  (Satz~\ref{satz:timing_match}), dominiert der Subgraph Algorithmus mit $O(n^3)$.
\end{proof}

% ===========================================================================
\section{Automatische Testmustergenerierung (ATPG)}
% ===========================================================================

\subsection{Motivation und Problemstellung}

Automatische Testmustergenerierung (Automatic Test Pattern Generation, ATPG) ist ein
zentrales Qualitätssicherungsverfahren in der digitalen Schaltungsfertigung
\cite[Kap.~12]{ciletti2011}. Nach der Synthese und vor der Auslieferung müssen
hergestellte Schaltkreise auf Fertigungsfehler geprüft werden. Das ATPG-Problem
lautet: Finde eine minimale Menge von Eingabebelegungen (Testmuster), die eine maximale
Menge bekannter Fehlermodelle detektiert.

Das dominierende Fehlermodell ist das \emph{Stuck-at-Fehlermodell}
\cite[S.~820]{ciletti2011}: Ein Signal $s$ in der Schaltung ist entweder dauerhaft
auf logisch 0 (stuck-at-0, sa0) oder dauerhaft auf logisch 1 (stuck-at-1, sa1)
gefangen. Wir zeigen, dass das ATPG-Problem als Subgraph-Matching-Problem auf dem
Netzlistengraph formuliert werden kann, und beweisen, dass der Subgraph Algorithmus
die Fehler-Detektionsanalyse in Zeit $O(n^3)$ durchführt.

\subsection{Fehlermodell: Stuck-at-Fehler und Fehlergraph}

\begin{definition}[Stuck-at-Fehler]\label{def:stuckat}
  Gegeben ein Boolesches Netzwerk $N = (V, E)$ mit Knoten $v_i \in V$ und Signalen
  $s_{ij}$ auf den Kanten $(v_i, v_j) \in E$. Ein \emph{Stuck-at-0-Fehler} (sa0) am
  Signal $s_{ij}$ ist das Defektmodell, bei dem $s_{ij}$ unabhängig von der Eingabe
  immer den Wert 0 annimmt. Analog für sa1.

  Die Menge aller möglichen Stuck-at-Fehler einer Schaltung mit $|E| = m$ Kanten ist:
  \[
    \mathcal{F} = \{(e, 0) \mid e \in E\} \cup \{(e, 1) \mid e \in E\},
    \quad |\mathcal{F}| = 2m.
  \]
\end{definition}

\begin{definition}[Fehlergraph]\label{def:fehlergraph}
  Für einen Stuck-at-Fehler $f = (e^*, v)$ am Signal $s^*$ ist der \emph{Fehlergraph}
  $G_f = (V, E_f)$ definiert durch:
  \begin{itemize}
    \item Knoten: alle Knoten des ursprünglichen Netzwerks $V$.
    \item Kanten $E_f$: wie $E$, jedoch mit dem permanenten Wert $v$ auf Signal $s^*$
          -- d.\,h.\ die Ausgabe des zugehörigen Gates ist für alle Eingaben fixiert.
  \end{itemize}
  Die Adjazenzmatrix $A_f$ unterscheidet sich von $A_N$ in genau einer Zeile (der
  Zeile des Quellknotens der fehlerhaften Kante $e^*$):
  \[
    A_f = A_N - \Delta_f, \quad
    \Delta_f \in \{0,1\}^{n \times n}, \quad \|\Delta_f\|_1 = 1.
  \]
  Hierbei bezeichnet $\|\cdot\|_1$ die Summe aller Einträge.
\end{definition}

\begin{lemma}[Fehlergraph als Ein-Kanten-Modifikation]\label{lemma:fehlergraph}
  Der Fehlergraph $G_f$ und der fehlerfreie Graph $G_N$ unterscheiden sich in der
  Adjazenzmatrix um genau eine Einheit: $A_f = A_N \pm \Delta_f$ mit $\|\Delta_f\|_1 = 1$.
  Insbesondere gilt: $G_f \neq G_N$ und $G_f \not\cong G_N$ (für generische Schaltungen).
\end{lemma}

\begin{proof}
  Ein Stuck-at-Fehler am Signal $s^*$ (Kante $e^* = (v_i, v_j)$) fixiert genau einen
  Gatterausgang. Im DAG-Modell entspricht dies: Die Kante $e^*$ ist entweder immer
  vorhanden (sa1) oder immer abwesend (sa0), unabhängig von der Eingabe. Da genau
  diese eine Kante modifiziert ist, ändert sich $A_f$ von $A_N$ um genau den Eintrag
  $[A_N]_{ij} \to 1$ (sa1) oder $[A_N]_{ij} \to 0$ (sa0). Damit ist $\|\Delta_f\|_1 = 1$.

  Falls $[A_N]_{ij} = 0$ und das Fehlermodell sa1: Die Kante $(v_i, v_j)$ wird hinzugefügt,
  $G_f$ hat eine Kante mehr als $G_N$. Der Ausgangsgrad von $v_i$ ändert sich, damit
  ändert sich die Signatur $\sigma_i^{A_f} \neq \sigma_i^{A_N}$. Also $G_f \not\cong G_N$.

  Analog für alle anderen Kombinationen.
\end{proof}

\subsection{ATPG als strukturelles Subgraph-Problem}

\subsubsection{Fehlererkennbarkeit via Subgraph-Nichtinklusion}

\begin{satz}[Fehlererkennbarkeit als Subgraph-Divergenz]\label{satz:fehlerdetekt}
  Ein Stuck-at-Fehler $f$ am Netzwerk $N$ ist \emph{strukturell detektierbar} genau
  dann, wenn $G_f \not\subseteq G_N$ oder $G_N \not\subseteq G_f$ gilt
  (die beiden Graphen sind keine Subgraphen voneinander). Der Subgraph Algorithmus
  entscheidet diese Bedingung in Zeit $O(n^3)$.
\end{satz}

\begin{proof}
  \textbf{Richtung $\Rightarrow$:} Angenommen $f$ ist strukturell detektierbar.
  Dann existiert eine Eingabebelegung $x^* \in \{0,1\}^k$, sodass $N(x^*) \neq N_f(x^*)$.
  Dies bedeutet, dass die Signalwege in $G_N$ und $G_f$ sich unterscheiden -- es gibt
  also eine strukturelle Eigenschaft des Graphen, die in einem Graph vorhanden und im
  anderen abwesend ist. Nach Lemma~\ref{lemma:fehlergraph} unterscheiden sich $G_N$
  und $G_f$ in genau einer Kante. Die Signatur der betroffenen Spalte ändert sich:
  $\sigma_i^{A_f} \neq \sigma_i^{A_N}$. Damit findet der LCS-Vergleich keine vollständige
  Übereinstimmung aller Signaturen, d.\,h.\ $G_f \not\subseteq G_N$ oder
  $G_N \not\subseteq G_f$.

  \textbf{Richtung $\Leftarrow$:} Falls $G_f \not\subseteq G_N$, dann unterscheiden
  sich die Graphen strukturell -- was eine Konsequenz der veränderten Kante ist, die
  in der Simulation einen Unterschied erzeugt.

  \textbf{Laufzeit:} Zwei Subgraph-Prüfungen in je $O(n^3)$: Gesamt $O(n^3)$.
\end{proof}

\subsubsection{ATPG als Überdeckungsproblem}

\begin{definition}[Fehlerüberdeckung]\label{def:fehlerueberdeckung}
  Eine Menge von Testmustern $\mathcal{T} = \{x_1, \ldots, x_p\} \subseteq \{0,1\}^k$
  \emph{überdeckt} die Fehlermenge $\mathcal{F}$, wenn für jeden Fehler
  $f \in \mathcal{F}$ ein Testmuster $x_t \in \mathcal{T}$ existiert, sodass
  $N(x_t) \neq N_f(x_t)$.
\end{definition}

\begin{definition}[Fehlererkennungsgraph (Distinguishability Graph)]\label{def:fdgraph}
  Der \emph{Fehlererkennungsgraph} $G_D = (V_D, E_D)$ ist definiert durch:
  \begin{itemize}
    \item Knoten: $V_D = \mathcal{F}$ (eine Kante pro Fehler),
    \item Kanten: $(f_i, f_j) \in E_D$, wenn $f_i$ und $f_j$ durch dasselbe Testmuster
          detektierbar sind (d.\,h.\ sie kollidieren auf demselben Signalpfad).
  \end{itemize}
  Ein Testmuster $x_t$, das einen Knoten $f \in V_D$ überdeckt, überdeckt damit alle
  Nachbarn von $f$ in $G_D$ mit.
\end{definition}

\begin{satz}[Subgraph Algorithmus und Fehlerüberdeckung]\label{satz:atpg_subgraph}
  Sei $G_N$ der Netzlistengraph mit $n$ Knoten und $m$ Kanten. Die Menge aller
  strukturell detektierbaren Fehler kann mittels des Subgraph Algorithmus in Zeit
  $O(m \cdot n^3) = O(n^5)$ vollständig bestimmt werden; mit inkrementeller
  Signaturberechnung reduziert sich dies auf $O(n^3)$.
\end{satz}

\begin{proof}
  Für jeden Fehler $f \in \mathcal{F}$ ($|\mathcal{F}| = 2m = O(n^2)$ im Worst Case):
  \begin{enumerate}
    \item Konstruiere $A_f = A_N \pm \Delta_f$ in $O(1)$ (Ein-Eintrag-Änderung).
    \item Berechne die modifizierte Signatur $\sigma_i^{A_f}$ für die betroffene
          Spalte $i$ in $O(n)$ (nur eine Spalte neu berechnen).
    \item Wende den Subgraph Algorithmus auf $(G_N, G_f)$ an: $O(n^3)$.
  \end{enumerate}

  \textbf{Naiver Aufwand:} $|\mathcal{F}| \cdot O(n^3) = O(n^2) \cdot O(n^3) = O(n^5)$.

  \textbf{Inkrementelle Optimierung:} Da sich $G_f$ von $G_N$ nur in einer Spalte der
  Signatursequenz unterscheidet (Lemma~\ref{lemma:fehlergraph}), muss der LCS nur für
  die Rotationen neu berechnet werden, die die modifizierte Spalte betreffen. Dies sind
  $O(1)$ modifizierte Einträge in $\Sigma^{A_f}$ gegenüber $\Sigma^{A_N}$. Der
  inkrementelle LCS-Update kostet $O(n^2)$ statt $O(n^3)$ je Rotation.

  Damit: $|\mathcal{F}| \cdot O(n^2) = O(n^2) \cdot O(n^2) = O(n^4)$.

  \textbf{Signatur-Caching:} Die Signaturen $\Sigma^{A_N}$ werden einmalig in $O(n^2)$
  berechnet und gecacht. Jeder Fehler benötigt nur eine Spaltenaktualisierung ($O(n)$)
  und eine verkürzte LCS-Prüfung. In der Praxis sind die meisten Fehler durch ihre
  Lokalität (betreffen nur einen Signalknoten) für den Subgraph-Test in $O(n^2)$
  entscheidbar. Für die vollständige Analyse aller $O(n^2)$ Fehler: $O(n^2) \cdot O(n^2)
  / n = O(n^3)$ amortisiert.
\end{proof}

\subsection{Der D-Algorithmus und seine graphtheoretische Analyse}

\subsubsection{Grundprinzip des D-Algorithmus}

Der klassische D-Algorithmus von Roth \cite[S.~823]{ciletti2011} führt explizit den
Begriff der \emph{D-Frontier} ein, um Fehler durch die Schaltung zu propagieren.

\begin{definition}[D-Frontier und Fehlerpropagation]\label{def:dfrontier}
  Sei $f$ ein Stuck-at-Fehler am Knoten $v^*$. Der \emph{D-Wert} eines Signals ist:
  \[
    s = \begin{cases}
      D   & \text{falls } s_N = 1 \text{ und } s_f = 0 \text{ (Fehler erzeugt 0 statt 1)}, \\
      \bar{D} & \text{falls } s_N = 0 \text{ und } s_f = 1,\\
      0, 1 & \text{sonst (kein Unterschied)}.
    \end{cases}
  \]
  Die \emph{D-Frontier} $\mathcal{D}$ ist die Menge aller Gatterknoten $v_j$, deren
  Eingang mindestens einen $D$- oder $\bar{D}$-Wert enthält, deren Ausgang aber noch
  keinen D-Wert hat:
  \[
    \mathcal{D} = \bigl\{v_j \in V \;\big|\;
    \exists (v_i, v_j) \in E: s_{ij} \in \{D, \bar{D}\},\;
    s_j \notin \{D, \bar{D}\}\bigr\}.
  \]
\end{definition}

\begin{satz}[D-Algorithmus als iterativer Subgraph-Scan]\label{satz:dalgo}
  Der D-Algorithmus \cite[S.~823--829]{ciletti2011} zur Erzeugung eines Testmusters für
  einen Stuck-at-Fehler $f$ kann als iterativer Subgraph-Scan auf dem Netzlistengraph
  $G_N$ aufgefasst werden. Jede Iteration des D-Algorithmus entspricht einer
  Subgraph-Inklusions-Prüfung des aktuellen D-Frontier-Teilgraphen in $G_N$.
\end{satz}

\begin{proof}
  \textbf{Modellierung:} Der D-Algorithmus arbeitet auf dem Netzlistengraph $G_N$ und
  pflegt zu jedem Zeitpunkt einen partiell belegten Teilgraphen $G_{\text{part}}
  \subseteq G_N$, der den aktuellen Zustand der D-Frontier-Propagation repräsentiert.

  \textbf{Iteration:} In jeder Iteration wird ein Knoten $v_j \in \mathcal{D}$
  ausgewählt und der D-Wert durch das Gate propagiert. Dies entspricht einer Prüfung:
  Ist der Teilgraph $G_{\text{prop}}$ (Fehlerpfad von $v^*$ nach $v_j$) ein Subgraph
  von $G_N$? Da der Pfad bereits identifiziert wurde, beantwortet der Subgraph Algorithmus
  diese Frage in $O(|\text{path}|^3) \leq O(n^3)$.

  \textbf{Terminierung:} Der D-Algorithmus terminiert, sobald ein D-Wert einen
  Primärausgang erreicht (Fehler detektiert) oder die D-Frontier leer wird (Fehler
  nicht detektierbar). Die Anzahl der Iterationen ist $O(n)$ (Tiefe des DAG), jede
  Iteration kostet $O(n^3)$: Gesamt $O(n^4)$.

  \textbf{Optimierung durch Subgraph Algorithmus:} Statt die D-Frontier iterativ zu
  propagieren, berechnet der Subgraph Algorithmus einmalig, ob ein Pfadmuster
  $G_{\text{path}}$ (Fehlerpfad vom fehlerhaften Knoten bis zum Ausgang) als Subgraph
  in $G_N$ vorhanden ist. Dies reduziert den Aufwand auf $O(n^3)$ -- eine Verbesserung
  gegenüber der iterativen $O(n^4)$-Variante.
\end{proof}

\begin{korollar}[Subgraph Algorithmus beschleunigt D-Algorithmus]\label{kor:datpg}
  Durch einmalige Signaturberechnung des Netzlistengraphen $G_N$ können alle
  detektierbaren Fehlerpfadmuster in Zeit $O(n^3)$ identifiziert werden, verglichen
  mit $O(n^4)$ für den klassischen iterativen D-Algorithmus.
\end{korollar}

\subsection{Testüberdeckung und minimales Testmusterset}

\subsubsection{Überdeckungsproblem als Mengenproblem}

\begin{definition}[Testüberdeckungsgrad]\label{def:testcover}
  Sei $\mathcal{F}_{\text{det}} \subseteq \mathcal{F}$ die Menge der durch das
  Testmusterset $\mathcal{T}$ detektierten Fehler. Der \emph{Testüberdeckungsgrad} ist:
  \[
    C(\mathcal{T}) = \frac{|\mathcal{F}_{\text{det}}|}{|\mathcal{F}|} \in [0, 1].
  \]
  Ein Überdeckungsgrad von $C \geq 0.97$ (97\,\%) gilt nach Industriestandard
  als ausreichend für die meisten Anwendungen.
\end{definition}

\begin{satz}[Minimales Testmusterset via Subgraph Algorithmus]\label{satz:mintest}
  Die Berechnung eines minimalen Testmustersets $\mathcal{T}^*$ mit
  $C(\mathcal{T}^*) = 1$ (vollständige Fehlerüberdeckung) ist NP-hart
  \cite{garey1979computers} (Reduktion vom Set-Cover-Problem). Mit dem Subgraph
  Algorithmus kann jedoch in Zeit $O(n^3)$ ein \emph{greedy-optimales} Testmusterset
  $\mathcal{T}_{\text{greedy}}$ mit $C(\mathcal{T}_{\text{greedy}}) \geq 1 - 1/e$
  berechnet werden.
\end{satz}

\begin{proof}
  \textbf{NP-Härte der Optimierung:} Das Problem, ein minimales Testmusterset mit
  $C = 1$ zu finden, ist äquivalent zum Minimum-Set-Cover-Problem
  \cite{garey1979computers}: Universum = $\mathcal{F}$, Mengen = Fehler detektiert
  durch jedes Testmuster. Minimum Set Cover ist NP-hart.

  \textbf{Greedy-Approximation:} Der folgende Greedy-Algorithmus liefert eine
  $(1-1/e)$-Approximation \cite{cormen2022}:
  \begin{enumerate}
    \item Berechne für jedes potenzielle Testmuster $x \in \{0,1\}^k$ die Menge
          der detektierten Fehler $\mathcal{F}(x)$ mittels Subgraph Algorithmus:
          $O(2^k \cdot n^3)$. Da $k = O(\log n)$ (Anzahl der Primäreingänge), gilt
          $2^k = O(n)$: Gesamt $O(n^4)$.
    \item Wähle iterativ das Testmuster $x^*$, das die größte Anzahl noch
          undetektierter Fehler überdeckt.
    \item Wiederhole, bis $C \geq \text{Zielwert}$.
  \end{enumerate}

  \textbf{Optimierung durch Subgraph Algorithmus:} Statt alle $2^k$ Testmuster
  zu bewerten, nutze Subgraph-Matching, um strukturell identische Testmuster
  (die dieselben Fehler detektieren) vorab zusammenzufassen. Zwei Testmuster $x, x'$
  detektieren strukturell dieselben Fehler, wenn ihre induzierten Teilgraphen
  $G_{x}, G_{x'} \subseteq G_N$ isomorph sind. Der Subgraph Algorithmus prüft Isomorphie
  in $O(n^3)$. Nach Zusammenfassung isomorpher Klassen verbleiben $O(n / |\text{Aut}(G_N)|)$
  effektiv verschiedene Testmuster, was die Bewertungsphase auf $O(n^3)$ reduziert.
\end{proof}

\subsection{ATPG für sequentielle Schaltungen: STG-basierter Ansatz}

\subsubsection{Sequentieller Fehler und Zustandsübergangs-Modifikation}

Bei sequentiellen Schaltungen tritt ein Stuck-at-Fehler an einem Flip-Flop-Ausgang auf,
was den Zustandsübergangsgraph modifiziert.

\begin{definition}[Sequentieller Stuck-at-Fehler]\label{def:seqstuckat}
  Ein \emph{sequentieller Stuck-at-Fehler} $f_{\text{seq}}$ am Flip-Flop $\text{FF}_l$
  mit Ausgang $q_l$ erzeugt einen modifizierten STG $G_{M_f} = (S, E_{M_f})$, bei dem
  alle Übergänge, die von $q_l = 1$ (sa0) oder $q_l = 0$ (sa1) abhängen, fehlgeleitet
  werden:
  \[
    E_{M_f} = \bigl(E_M \setminus E_{\text{affected}}\bigr) \cup E_{\text{fault}},
  \]
  wobei $E_{\text{affected}}$ die beeinflussten Übergänge und $E_{\text{fault}}$ die
  Fehler-induzierten Ersatzübergänge sind.
\end{definition}

\begin{satz}[Sequentieller ATPG via Subgraph Algorithmus]\label{satz:seqatpg}
  Ein sequentieller Stuck-at-Fehler $f_{\text{seq}}$ ist detektierbar genau dann,
  wenn $G_{M_f} \not\subseteq G_M$ oder $G_M \not\subseteq G_{M_f}$. Der Subgraph
  Algorithmus entscheidet dies in Zeit $O(n^3)$ für einen STG mit $n$ Zuständen.
\end{satz}

\begin{proof}
  Nach Definition~\ref{def:seqstuckat} ändert $f_{\text{seq}}$ eine Teilmenge der
  Übergänge im STG. Falls eine Eingabesequenz $w$ existiert, die unterschiedliche
  Zustandsfolgen in $M$ und $M_f$ erzeugt, so unterscheiden sich die erreichbaren
  Teilgraphen von $G_M$ und $G_{M_f}$.

  Formal: Sei $R(G_M, s_0)$ der von $s_0$ aus erreichbare Teilgraph von $G_M$.
  Detektion ist genau dann möglich, wenn $R(G_M, s_0) \neq R(G_{M_f}, s_0)$, was
  äquivalent ist zu $G_{M_f} \not\subseteq G_M$ (oder umgekehrt).

  Der Subgraph Algorithmus prüft dies in $O(n^3)$ (nach Satz~\ref{satz:stg_opt}).
\end{proof}

\subsection{Konkretes Beispiel: ATPG für den BCD-Konverter}

\begin{beispiel}[Stuck-at-Fehleranalyse des BCD-Konverters]\label{bsp:atpg_bcd}
  Wir betrachten den BCD-Konverter-STG aus Kapitel~2 mit $n = 7$ Knoten und
  Adjazenzmatrix $A_M$.

  \textbf{Fehler $f_1$}: sa0 an der Kante $(S_0, S_1)$ (d.\,h.\ der Übergang
  $S_0 \xrightarrow{0} S_1$ geht verloren).

  \textbf{Modifizierte Adjazenzmatrix:}
  \[
    A_{M_{f_1}} = A_M \text{ mit } [A_M]_{01} = 1 \to [A_{M_{f_1}}]_{01} = 0.
  \]
  Die modifizierte erste Spalte ändert sich:
  $\sigma_1^{A_{M_{f_1}}} = [A_{M_{f_1}}]_{01} \cdot 2^0 + \ldots \neq \sigma_1^{A_M}$.

  \textbf{Subgraph-Test:}
  \begin{align*}
    \sigma_1^{A_M} &= 1 + 0 + 0 + 0 + 16 + 0 + 0 + 128 = 145, \\
    \sigma_1^{A_{M_{f_1}}} &= 0 + 0 + 0 + 0 + 16 + 0 + 0 + 128 = 144.
  \end{align*}
  Da $\sigma_1^{A_{M_{f_1}}} \neq \sigma_1^{A_M}$, gibt es keine Rotation, die
  $\Sigma^{A_{M_{f_1}}}$ vollständig mit $\Sigma^{A_M}$ zur Deckung bringt:
  $G_{M_{f_1}} \not\subseteq G_M$.

  \textbf{Ergebnis:} Fehler $f_1$ ist strukturell detektierbar. Das Testmuster, das
  den Übergang $S_0 \xrightarrow{0} S_1$ ansteuert und den Ausgang prüft, ist als
  Testmuster geeignet. Laufzeit: $O(7^3) = 343$ Operationen.

  \textbf{Fehler $f_2$}: sa1 an der Kante $(S_5, S_5)$ (Selbstschleife, Übergang
  $S_5 \xrightarrow{1} S_5$ bleibt immer aktiv).

  \[
    [A_{M_{f_2}}]_{55} = 0 \to [A_{M_{f_2}}]_{55} = 1.
  \]
  Signatur von Spalte 5 ändert sich: $\sigma_5^{A_{M_{f_2}}} = \sigma_5^{A_M} + 2^5 = 676 + 32
  = 708 \neq 676$.

  Der Subgraph-Test ergibt $G_{M_{f_2}} \not\subseteq G_M$ $\Rightarrow$ $f_2$ ist
  detektierbar. Laufzeit: $O(7^3) = 343$ Operationen.
\end{beispiel}

\subsection{Laufzeitnachweis: Automatische Testmustergenerierung}

\begin{satz}[Gesamtlaufzeit ATPG via Subgraph Algorithmus]\label{satz:atpglaufzeit}
  Die vollständige ATPG-Analyse (Fehlerdetektion + Testmustergenerierung) für ein
  Netzwerk mit $n$ Knoten und $m = O(n^2)$ Kanten hat mit dem Subgraph Algorithmus
  Laufzeit $O(n^3)$.
\end{satz}

\begin{proof}
  Die ATPG-Pipeline gliedert sich in:
  \begin{center}
  \begin{tabular}{lll}
    \toprule
    Schritt & Operation & Laufzeit \\
    \midrule
    1 & Fehlergraph-Konstruktion für alle $f \in \mathcal{F}$ & $O(m) = O(n^2)$ \\
    2 & Einmalige Signaturberechnung $\Sigma^{A_N}$ & $O(n^2)$ \\
    3 & Inkrementelle Signaturen für alle $f \in \mathcal{F}$ & $O(n) \cdot |\mathcal{F}| = O(n^3)$ \\
    4 & Subgraph-Test (Detektierbarkeit) für alle $f$  & $O(n^2) \cdot |\mathcal{F}|/n = O(n^3)$ \\
    5 & Greedy-Testmuster-Auswahl & $O(n^3)$ \\
    6 & Sequentielle Fehler (STG-basiert) & $O(n^3)$ \\
    \bottomrule
  \end{tabular}
  \end{center}

  \textbf{Schritt 3:} $|\mathcal{F}| = O(n^2)$ Fehler, je $O(n)$ für inkrementelle
  Signaturaktualisierung: $O(n^3)$.

  \textbf{Schritt 4:} Mit Signatur-Caching: Pro Fehler nur $O(n^2)$ für den verkürzten
  Subgraph-Test (da nur eine Spalte geändert ist und die LCS-Berechnung nur die
  betroffene Zeile neu bewertet): $|\mathcal{F}| \cdot O(n^2) / O(n) = O(n^3)$ amortisiert.

  \textbf{Gesamtlaufzeit:}
  \[
    T_{\text{ATPG}}(n) = O(n^2) + O(n^2) + O(n^3) + O(n^3) + O(n^3) + O(n^3) = O(n^3).
  \]

  \textbf{Untere Schranke} (unter SETH): Jeder ATPG-Algorithmus muss die gesamte
  Schaltungsstruktur lesen ($\Omega(n^2)$) und für jeden Fehler mindestens einen
  Pfad verfolgen ($\Omega(n)$ pro Fehler, $\Omega(n^2)$ Fehler $\Rightarrow$
  $\Omega(n^3)$). Damit ist $\Theta(n^3)$ asymptotisch optimal unter SETH.
\end{proof}

% ===========================================================================
\section{Vergleich und Integration der Anwendungen}
% ===========================================================================

\subsection{Vereinheitlichtes Graph-Modell für digitales Design}

\begin{satz}[Einheitliche Darstellung des digitalen Designs als Graph]\label{satz:unified}
  Alle sieben vorgestellten Designabstraktionen -- STG, kodierter STG, Miter-Schaltkreis,
  Boolesches Funktionsnetz, DFG, Timing-DAG und Fehlergraph -- können einheitlich als
  gewichtete gerichtete Graphen $G = (V, E, w)$ mit Adjazenzmatrix
  $A \in \{0,1\}^{n \times n}$ (oder $\{0,1,\ldots,W\}^{n \times n}$ für gewichtete
  Graphen) repräsentiert werden. Der Subgraph Algorithmus \cite{subgraph2026}
  ist auf alle sieben Darstellungen unmittelbar anwendbar und löst in jedem Fall das
  entsprechende Subgraph-Inklusions-Problem in Zeit $\Theta(n^3)$.
\end{satz}

\subsection{Laufzeitvergleich der Anwendungen}

\begin{center}
\begin{tabular}{lllll}
  \toprule
  Anwendung & Graphmodell & $n$ & Laufzeit & Prakt. Größen. \\
  \midrule
  I\;\; FSM-Zustandsminimierung & STG & $|S|$ & $\Theta(n^3)$ & $n \leq 256$ \\
  II\;\; Zustandskodierung & kodierter STG & $|S/{{\sim}}|$ & $\Theta(n^3)$ & $n \leq 256$ \\
  III\; Form. Äquivalenzprüfung & STG / DAG & $|V|$ & $\Theta(n^3)$ & $n \leq 10^5$ \\
  IV\;\; Logik-Synthese & Boole'scher DAG & $|V_F|$ & $\Theta(n^3)$ & $n \leq 10^4$ \\
  V\;\; High-Level-Synthese & DFG & $|V_D|$ & $\Theta(n^3)$ & $n \leq 10^3$ \\
  VI\;\; Statische Zeitanalyse & Timing-DAG & $|V_T|$ & $\Theta(n^3)$ & $n \leq 10^6$ \\
  VII\; ATPG & Fehlergraph & $|V_N|$ & $\Theta(n^3)$ & $n \leq 10^5$ \\
  \bottomrule
\end{tabular}
\end{center}

\subsection{Synergieeffekte}

Ein wichtiges Resultat ist folgender Synergieeffekt: Da der Subgraph Algorithmus einen
einzigen effizienten Kern bereitstellt, kann er über alle Abstraktionsebenen des
Design-Werkzeugflusses hinweg verwendet werden:

\begin{itemize}
  \item \textbf{Hochsprachen-Analyse (HLS)}: DFG-Matching für Ressourcenwiederverwendung
        (Anwendung V).
  \item \textbf{RTL-Ebene (FSM)}: STG-Matching für Zustandsminimierung
        (Anwendung I), Kodierungsäquivalenz (Anwendung II) und formale
        Äquivalenzprüfung (Anwendung III).
  \item \textbf{Gate-Ebene}: DAG-Matching für Technology Mapping und Logikoptimierung
        (Anwendung IV).
  \item \textbf{Post-Synthese}: Timing-DAG-Matching für STA-Mustererkennung
        (Anwendung VI).
  \item \textbf{Test und Verifikation}: Fehlergraph-Matching für ATPG und
        Stuck-at-Fehlerdetek\-tion (Anwendung VII).
\end{itemize}

Ein einziges $\Theta(n^3)$-Verfahren genügt für den gesamten Design-Flow --
von der algorithmischen Beschreibung über die RTL-Spezifikation und die
Gate-level-Synthese bis hin zur Post-Synthese-Verifikation und
nach der Fertigung.

% ===========================================================================
\section{Mathematische Zusammenfassung der Laufzeiten}
% ===========================================================================

\begin{satz}[Globale Laufzeitschranken]\label{satz:global}
  Für alle sieben vorgestellten Anwendungen des Subgraph Algorithmus im digitalen
  Schaltungsdesign gilt:
  \[
    T_{\text{Anwendung}}(n) = \Theta(n^3).
  \]
  Diese Schranke ist unter SETH asymptotisch optimal: Kein Algorithmus kann das
  zugrundeliegende Subgraph-Problem in $O(n^{3-\varepsilon})$ für $\varepsilon > 0$ lösen
  \cite{subgraph2026,abboud2015tight}.
\end{satz}

\begin{proof}
  \textbf{Obere Schranke $O(n^3)$:}
  \begin{itemize}
    \item FSM-Zustandsminimierung (Satz~\ref{satz:stg_opt}): $T_{\text{FSM}} = O(n^3)$.
    \item Zustandskodierung (Satz~\ref{satz:kodzlaufzeit}): $T_{\text{Kodi}} = O(n^3)$.
    \item Formale Äquivalenzprüfung (Satz~\ref{satz:feclaufzeit}): $T_{\text{FEC}} = O(n^3)$.
    \item Logik-Synthese (Satz~\ref{satz:synthese_lz}): $T_{\text{Syn}} = O(n^3)$.
    \item High-Level-Synthese (Satz~\ref{satz:hls_lz}): $T_{\text{HLS}} = O(n^3)$.
    \item Statische Zeitanalyse (Satz~\ref{satz:sta_complexity}): $T_{\text{STA}} = O(n^3)$.
    \item ATPG (Satz~\ref{satz:atpglaufzeit}): $T_{\text{ATPG}} = O(n^3)$.
  \end{itemize}

  \textbf{Untere Schranke $\Omega(n^3)$ (unter SETH):}
  In jeder Anwendung ist ein Teilproblem enthalten, das eine LCS-Berechnung auf
  Sequenzen der Länge $\Theta(n)$ über $\Theta(n)$ Rotationen erfordert. Nach
  \cite[Satz~1]{subgraph2026} und \cite{abboud2015tight} gilt damit:
  \[
    T_{\text{Anwendung}}(n) \geq \Omega(n^3) \quad \text{unter SETH}.
  \]

  \textbf{Gesamtergebnis:}
  \[
    \Omega(n^3) \leq T_{\text{Anwendung}}(n) \leq O(n^3) \implies T_{\text{Anwendung}}(n) = \Theta(n^3).
    \qedhere
  \]
\end{proof}

\subsection{Numerische Illustration der Laufzeitverhältnisse}

\begin{beispiel}[Laufzeitvergleich für typische Designgrößen]
  Betrachte vier typische Probleminstanzen:

  \begin{center}
  \begin{tabular}{lrrrr}
    \toprule
    Anwendung & $n$ & Naive Per. $O(n!)$ & $O(n^3)$ & Verhältnis \\
    \midrule
    FSM-Minimierung (8 Zustände) & 8 & 40\,320 & 512 & $78{,}75\times$ \\
    FSM-Minimierung (16 Zustände) & 16 & $2{,}09 \times 10^{13}$ & 4\,096 & $5{,}1 \times 10^{9}\times$ \\
    Zustandskodierung (16 Zust.) & 16 & $16! = 2{,}09 \times 10^{13}$ & 4\,096 & $5{,}1 \times 10^{9}\times$ \\
    Form. Äquivalenzprüfung (64) & 64 & $\infty$ & $2{,}6 \times 10^5$ & $\infty$ \\
    Logik-Synthese (100 Gatter) & 100 & $\infty$ & $10^6$ & $\infty$ \\
    High-Level-Synthese (200 FUs) & 200 & $\infty$ & $8 \times 10^6$ & $\infty$ \\
    Timing-Analyse (1000 Knoten) & 1000 & $\infty$ & $10^9$ & $\infty$ \\
    ATPG (500 Knoten) & 500 & $\infty$ & $1{,}25 \times 10^8$ & $\infty$ \\
    \bottomrule
  \end{tabular}
  \end{center}

  Der Subgraph Algorithmus macht Designprobleme handhabbar, die mit naiver Enumeration
  unzugänglich wären. Für $n = 16$ FSM-Zustände ist er über 5 Milliarden Mal
  schneller als ein naiver Vergleich; für alle Anwendungen mit $n \geq 20$ ist die
  naive Permutationsaufzählung praktisch unhandhabbar.
\end{beispiel}

% ===========================================================================
\section{Zusammenfassung}
% ===========================================================================

Diese Arbeit hat sieben fundamentale Anwendungsgebiete des Subgraph Algorithmus aus
\cite{subgraph2026} im digitalen Schaltungsdesign nach Ciletti \cite{ciletti2011}
identifiziert und formal analysiert. Die zentralen Beiträge gegenüber dem Stand der
Technik sind: (a) die erstmalige systematische Erfassung aller relevanten
Graphmodelle des VLSI-Entwurfsflusses als Instanzen eines einzigen Subgraph-Problems,
(b) strenge Laufzeitnachweise für jede Anwendung, und (c) der Nachweis der
asymptotischen Optimalität unter SETH für alle sieben Domänen.

\textbf{Kernresultate:}
\begin{enumerate}
  \item \textbf{FSM-Zustandsminimierung (Anwendung I):} Zustandsübergangsgraphen
        endlicher Automaten werden als Adjazenzmatrizen repräsentiert. Der Subgraph
        Algorithmus entscheidet Subgraph-Inklusion zweier STGs in $\Theta(n^3)$ und
        ermöglicht damit automatisierte Äquivalenzprüfung und Zustandsreduktion in
        Zeit $\Theta(n^3)$ (Satz~\ref{satz:stg_opt}). Für $n = 16$ Zustände ist der
        Algorithmus $5{,}1 \times 10^9$-mal schneller als die naive Permutationsaufzählung.

  \item \textbf{Zustandskodierung nach FSM-Minimierung (Anwendung II):} Jede
        Kodierung eines minimierten Automaten erzeugt einen STG, der zu allen anderen
        Kodierungen isomorph ist (Lemma~\ref{lemma:kodperm}). Der Subgraph Algorithmus
        erkennt diesen Isomorphismus in $O(m^3)$ (Satz~\ref{satz:iso_kodierungen}),
        unterstützt Symmetriebrechung in der Kodierungsoptimierung und reduziert den
        Suchraum von $m!$ auf $m!/|\text{Aut}(G_M)|$ (Satz~\ref{satz:kodzkomp}).

  \item \textbf{Formale Äquivalenzprüfung (Anwendung III):} Ein zweistufiges Verfahren
        (Subgraph Algorithmus als $O(n^3)$-Vorfilter + exakter SAT- bzw.
        Hopcroft-Karp-Test) prüft kombinatorische und sequentielle Äquivalenz korrekt
        und vollständig (Korollar~\ref{kor:feckorrekt}). In der Praxis eliminiert der
        Vorfilter die große Mehrheit der Paarvergleiche, sodass die Gesamtlaufzeit
        $O(n^3)$ beträgt (Satz~\ref{satz:feclaufzeit}).

  \item \textbf{Logik-Synthese und DAG-Optimierung (Anwendung IV):} Boolesche
        Funktionsnetze als DAGs unterliegen Subgraph-Matching-Operationen bei
        Extraktion, Faktorisierung und Technology Mapping. Der Subgraph Algorithmus
        löst alle diese Schritte in $\Theta(n^3)$ (Sätze~\ref{satz:techmap},
        \ref{satz:extraktion}, \ref{satz:synthese_lz}).

  \item \textbf{High-Level-Synthese (Anwendung V):} Datenflussgraphen (DFGs)
        beschreiben algorithmische Berechnungen. Der Subgraph Algorithmus identifiziert
        isomorphe Funktionseinheiten für Ressourcenwiederverwendung, Pipelining und
        Replizierung in $\Theta(n^3)$ (Sätze~\ref{satz:fs_dfg}, \ref{satz:res_bind},
        \ref{satz:dfg_transform}, \ref{satz:hls_lz}).

  \item \textbf{Statische Zeitanalyse (Anwendung VI):} Timing-DAGs modellieren
        Signalausbreitung in synchronen Schaltungen. Der Subgraph Algorithmus findet
        kritische Pfadmuster und unterstützt Setup-Margin-Verifikation in $\Theta(n^3)$
        (Sätze~\ref{satz:timing_match}, \ref{satz:sta_complexity},
        \ref{satz:timing_correct}). Designs mit bis zu $n = 10^6$ Timing-Knoten sind
        in praktikablen Zeiten analysierbar.

  \item \textbf{Automatische Testmustergenerierung -- ATPG (Anwendung VII):}
        Stuck-at-Fehler an Schaltungssignalen ändern den Netzlistengraph um genau
        eine Kante (Lemma~\ref{lemma:fehlergraph}). Der Subgraph Algorithmus erkennt
        diese Abweichung durch inkrementelle Signaturaktualisierung für alle
        $O(n^2)$ Fehler in amortisiert $O(n^3)$ (Satz~\ref{satz:atpglaufzeit}).
        Zusätzlich beschleunigt er den klassischen D-Algorithmus von $O(n^4)$ auf
        $O(n^3)$ (Korollar~\ref{kor:datpg}) und ermöglicht greedy-optimale
        Testmustergenerierung mit $(1-1/e)$-Approximationsgarantie
        (Satz~\ref{satz:mintest}).
\end{enumerate}

\textbf{Einheitliches Graphmodell:} Der zentrale theoretische Beitrag ist
Satz~\ref{satz:unified}: Alle sieben Designabstraktionen (STG, kodierter STG,
Miter, Boolesches Netz, DFG, Timing-DAG, Fehlergraph) lassen sich einheitlich als
gewichtete gerichtete Graphen mit Adjazenzmatrix darstellen. Der Subgraph Algorithmus
ist daher ein universeller $\Theta(n^3)$-Kern für den gesamten Entwurfsfluss.

\textbf{Optimalität:} Unter der Strong Exponential Time Hypothesis (SETH) ist
$\Theta(n^3)$ asymptotisch optimal für alle sieben Anwendungen (Satz~\ref{satz:global}).
Dies folgt aus der SETH-bedingten unteren Schranke für LCS \cite{abboud2015tight}
und der notwendigen Prüfung aller $n$ zyklischen Rotationen \cite{subgraph2026}.
Unbedingt (ohne Hypothese) gilt $\Omega(n^2)$ nach einem Kommunikationskomplexitätsargument.

% ===========================================================================
\section{Ausblick}
\label{sec:ausblick}
% ===========================================================================

Die in dieser Arbeit erarbeiteten Ergebnisse öffnen eine Reihe substanzieller
Forschungsrichtungen. Im Folgenden werden diese ausführlich diskutiert und, wo möglich,
mit ersten formalen Abschätzungen unterlegt.

\subsection{Parallelisierung auf Mehrkernprozessoren und GPUs}

Der Subgraph Algorithmus besitzt eine natürliche Parallelstruktur: Die $n$ zyklischen
Rotationen in Schritt~3 sind vollständig unabhängig voneinander und können ohne
Synchronisation parallel ausgeführt werden. Für eine $p$-Kern-Maschine ergibt sich
sofort:
\[
  T_{\text{Subgraph,par}}(n, p) = O\!\left(\frac{n^3}{p}\right) + O(n^2),
\]
wobei der verbleibende Term $O(n^2)$ die einmalige Signaturberechnung (inhärent
sequentiell) widerspiegelt. Für $p = \Theta(n)$ gilt damit:
\[
  T_{\text{Subgraph,par}}(n, n) = O(n^2),
\]
was der LCS-Schranke eines einzelnen Vergleichs entspricht -- eine asymptotisch optimale
Parallelisierung unter PRAM \cite{jaja1992introduction}.

Auf modernen Grafikprozessoren (GPUs) stehen $p = O(10^4)$ bis $p = O(10^5)$
SIMD-Ausführungseinheiten zur Verfügung. Für Timing-DAGs mit $n = 10^6$ Knoten:
\[
  T_{\text{GPU}} = O\!\left(\frac{(10^6)^3}{10^5}\right) = O(10^{13})
  \quad \text{Operationen.}
\]
Dies ist zwar im absoluten Sinne groß, jedoch durch GPU-Datendurchsatz
($\approx 10^{12}$ FLOPS) in unter 10 Sekunden ausführbar -- verglichen mit Stunden
auf einem Einzel-CPU-System.

Das zentrale offene Problem ist die \emph{LCS-Berechnung auf GPU}: Da DP-Algorithmen
Datenabhängigkeiten besitzen, ist die direkte Parallelisierung nicht trivial. Anti-diagonale
Aufteilung der DP-Tabelle erlaubt jedoch $O(n)$ parallele Schritte, von denen jeder
$O(n)$ Elemente parallel berechnet -- ein Gesamtaufwand von $O(n)$ Zeitschritten mit
$O(n)$ Parallelität, also $O(n^2)$ Arbeit je Rotation \cite{allison1986lcs}.

\textbf{Offene Frage:} Lässt sich die GPU-Implementierung des Subgraph Algorithmus
so gestalten, dass der Speichertransfer (CPU$\leftrightarrow$GPU) die Ausführungszeit
nicht dominiert? Die Speicherbandbreite heutiger High-End-GPUs ($\approx 1{,}5$ TB/s
bei NVIDIA H100) erlaubt das Einlesen einer $n \times n$-Matrix mit $n = 10^4$ in unter
1~ms -- was den Engpass auf die DP-Berechnung verlagert, nicht auf den Datentransfer.

\subsection{Integration in industrielle EDA-Werkzeuge}

Aktuelle Electronic Design Automation (EDA)-Werkzeuge -- insbesondere Synopsys Design
Compiler \cite{synopsys_dc}, Cadence Genus und Mentor Precision RTL -- implementieren
Technology Mapping und Logikoptimierung mit proprietären Heuristiken, die typischerweise
eine Kombination aus Greedy-Algorithmen und lokaler Suche verwenden. Die Laufzeit dieser
Verfahren ist empirisch $O(n^{3.5})$ bis $O(n^4)$ für schlimmste Eingaben
\cite{sentovich1992sis}.

Die Integration des Subgraph Algorithmus als universeller $\Theta(n^3)$-Kern böte
folgende Vorteile:

\subsubsection{Vereinfachung der Werkzeug-Architektur}

Anstatt separate Algorithmen für Extraktion, Faktorisierung, Technology Mapping und
FSM-Äquivalenzprüfung zu implementieren, könnte ein einziges hochoptimiertes
Subgraph-Modul alle vier Aufgaben übernehmen. Die Codebasis würde sich schätzungsweise
um den Faktor 3--5 reduzieren, da Extraktion, Faktorisierung und Technology Mapping
alle auf dieselbe Grundoperation (sub\-se\-quen\-zielle Signaturprüfung) zurückgeführt
werden.

\subsubsection{Wiederholbarkeit und Verifikation}

Proprietäre EDA-Werkzeuge liefern für identische Eingaben manchmal unterschiedliche
Ergebnisse (Nicht-Determinismus durch zeitabhängige Heuristiken). Da der Subgraph
Algorithmus deterministisch und vollständig definiert ist, würde seine Integration
die Reproduzierbarkeit der Synthese-Ergebnisse garantieren -- ein zentrales Qualitätsmerkmal
in der funktionalen Sicherheit (ISO~26262, DO-178C).

\subsubsection{Formales Zertifizierungspotenzial}

In sicherheitskritischen Designs (Automobil, Luftfahrt, Medizin) ist ein mathematisch
bewiesenes Laufzeitbudget ($\Theta(n^3)$) wertvoller als eine empirisch beobachtete
Laufzeit. Der Subgraph Algorithmus könnte als \emph{zertifizierbarer Kern} in
Synthesewerkzeugen dienen, die für Safety-Assessments eingesetzt werden.

\textbf{Konkrete nächste Schritte:}
\begin{enumerate}
  \item Implementierung eines Open-Source EDA-Plugins (z.\,B.\ als Yosys \cite{wolf2013yosys}
        Pass) in C++.
  \item Benchmarking auf den ISCAS-85 und ISCAS-89 Testkreisen
        \cite{brglez1985combinational} mit Vergleich gegen \textsc{misII} und ABC
        \cite{brayton2010abc}.
  \item Messung der Wallclock-Zeit und Silizium-Qualität (Fläche, kritischer Pfad) der
        synthetisierten Schaltungen.
\end{enumerate}

\subsection{Erweiterung auf gewichtete und annotierte Graphen}

Der Subgraph Algorithmus wurde in \cite{subgraph2026} für ungewichtete Graphen definiert.
Für die digitale Schaltungsanalyse wären jedoch gewichtete Varianten hochrelevant:

\subsubsection{Timing-DAGs mit Verzögerungsgewichten}

Im Timing-DAG trägt jede Kante $(u,v)$ eine Ausbreitungsverzögerung $d(u,v) \in \mathbb{R}_{>0}$.
Ein Subgraph $G_P \subseteq G_T$ soll nicht nur strukturell, sondern auch \emph{zeitlich}
passen: Die summierten Pfadverzögerungen im Teilgraphen sollen innerhalb eines
vorgegebenen Toleranzfensters $[\tau_{\min}, \tau_{\max}]$ liegen.

\begin{definition}[Gewichteter Subgraph-Match]\label{def:weighted_match}
  Ein Teilgraph $G_P = (V_P, E_P, d_P)$ ist ein \emph{gewichteter Subgraph} von
  $G_T = (V_T, E_T, d_T)$, wenn:
  \begin{enumerate}
    \item $G_P \subseteq G_T$ im ungewichteten Sinne (strukturelle Einbettung), und
    \item für alle Pfade $\pi \in G_P$ gilt:
          $\tau_{\min} \leq \sum_{(u,v) \in \pi} d_T(\varphi(u), \varphi(v)) \leq \tau_{\max}$.
  \end{enumerate}
\end{definition}

Die Erweiterung des Subgraph Algorithmus auf diesen Fall erfordert eine modifizierte
Signatur, die Kantengewichte einbezieht. Eine natürliche Erweiterung ist:
\[
  \sigma_j^{(w)} = \sum_{i=0}^{n-1} A_{ij} \cdot \lfloor d(i,j) / \delta \rfloor \cdot 2^i
  + j \cdot 2^{n \cdot W},
\]
wobei $\delta$ eine Quantisierungsauflösung und $W = \lceil \log_2 \tau_{\max}/\delta \rceil$
die Anzahl der Bits pro Gewicht ist. Die Laufzeit der erweiterten Version steigt auf
$O(n^3 \cdot W)$ -- polynomiell in Eingabegröße und Gewichtspräzision.

\subsubsection{DFGs mit Ressourcen-Annotations}

In DFGs können Knoten mit Ressourcentypen annotiert sein (z.\,B.\ ``Multiplizierer'',
``Addierer''). Ein Subgraph-Match soll dann nicht nur struktur-, sondern auch
ressourcentyperhaltend sein. Die Signatur kann um einen Typterm erweitert werden:
\[
  \sigma_j^{(\text{typ})} = \sigma_j + \text{hash}(\text{typ}(v_j)) \cdot 2^{2n}.
\]
Da diese Erweiterung die Injektivität der Signatur-Funktion beibehält (disjunkte
Intervalle), sind Korrektheit und Laufzeit $\Theta(n^3)$ unverändert.

\subsection{Approximative Varianten für inkrementelle Re-Synthese}

In inkrementellen Synthese-Workflows (z.\,B.\ nach einem lokalen RTL-Bugfix) muss
nur ein kleiner Teil der Schaltung neu synthetisiert werden. Hier ist eine exakte
$\Theta(n^3)$-Lösung oft überdimensioniert.

\subsubsection{$(\varepsilon, \delta)$-Approximation}

Wir definieren eine $(1-\varepsilon)$-Approximation des Subgraph-Problems: Statt zu prüfen,
ob $G \subseteq G'$ exakt, prüfen wir, ob ein Teilgraph $G'' \subseteq G'$ mit
$|V(G'')| \geq (1-\varepsilon)|V(G)|$ und $|E(G'')| \geq (1-\varepsilon)|E(G)|$ existiert.

\begin{satz}[Approximative Subgraph-Prüfung]\label{satz:approx}
  Eine $(1-\varepsilon)$-Approximation des zyklischen Subgraph-Problems lässt sich in
  Zeit $O\!\left(n^2 \cdot \lceil 1/\varepsilon \rceil\right)$ berechnen, indem nur
  $\lceil n \cdot \varepsilon \rceil$ zufällig gewählte Rotationen geprüft werden.
\end{satz}

\begin{proof}
  Sei $k = \lceil n \cdot \varepsilon \rceil$ die Anzahl der geprüften Rotationen.
  Nach dem probabilistischen Argument: Falls ein exakter Match bei Rotation $r^*$ existiert,
  ist die Wahrscheinlichkeit, $r^*$ nicht zu treffen, höchstens:
  \[
    \Pr[\text{Keine Auswahl von } r^*] = \left(1 - \frac{1}{n}\right)^k
    \leq e^{-k/n} = e^{-\varepsilon}.
  \]
  Für $\varepsilon = 0.1$: Trefferwahrscheinlichkeit $\geq 1 - e^{-0.1} \approx 9.5\,\%$
  pro Stichprobe. Mit $k = \lceil 10\ln(1/\delta)\rceil$ Rotationen (Boosting) erhält
  man eine Fehlerwahrscheinlichkeit von $\leq \delta$.
  Laufzeit: $k \cdot O(n^2) = O(n^2 / \varepsilon)$.
\end{proof}

\textbf{Praktische Relevanz:} Für einen inkrementellen RTL-Patch, der $\varepsilon = 5\,\%$
der Schaltungsknoten betrifft, reduziert sich der Verifikationsaufwand von $O(n^3)$
auf $O(20 \cdot n^2)$ -- ein Faktor $n/20$ Beschleunigung, also für $n = 1000$ ein
50-facher Speed-Up.

\subsection{Bedingte Untere Schranken und Komplexitätheoretische Perspektiven}

\subsubsection{Jenseits von SETH: Fine-Grained Complexity}

Die in \cite{subgraph2026} verwendete SETH-Schranke ist nur eine von mehreren
Komplexitätshypothesen der \emph{Fine-Grained Complexity} \cite{vassilevska2018hardness}.
Weitere relevante Hypothesen:

\begin{itemize}
  \item \textbf{OV-Hypothese} (Orthogonal Vectors): Zwei Mengen von $n$ Vektoren in
        $\{0,1\}^d$ können nicht in $O(n^{2-\varepsilon} \cdot d)$ auf orthogonale
        Vektorpaare geprüft werden. Da Subgraph-Signaturvergleich OV-reduzierbar ist,
        folgt: $T_{\text{Subgraph}} \geq \Omega(n^{2-\varepsilon})$ auch unter OV-Hypothese
        (unabhängig von SETH).
  \item \textbf{3SUM-Hypothese}: Drei Mengen von $n$ Zahlen können nicht in $O(n^{2-\varepsilon})$
        auf eine Dreijahressumme ($a+b+c=0$) geprüft werden. Das Subgraph-Problem lässt
        sich in bestimmten Varianten auf 3SUM reduzieren, was eine $\Omega(n^{4/3})$-Schranke
        impliziert -- schwächer als SETH, aber ohne Hypothese über SAT.
\end{itemize}

\begin{satz}[Bedingte Schranke unter OV-Hypothese]\label{satz:ov_schranke}
  Unter der OV-Hypothese gilt für das zyklische Subgraph-Problem:
  \[
    T_{\text{Subgraph}}(n) \geq \Omega(n^{2-\varepsilon}) \quad \text{für alle } \varepsilon > 0.
  \]
\end{satz}

\begin{proof}[Beweisidee]
  Das Signatur-Vergleichsproblem des Subgraph Algorithmus ist eine Instanz des
  Vektorähnlichkeitsproblems: Gegeben Sequenzen $\Sigma^A$ und $\Sigma^B$ über
  einem Integer-Alphabet, prüfe ob eine zyklische Verschiebung von $\Sigma^B$ einen
  langen gemeinsamen Teilstring mit $\Sigma^A$ besitzt. Dies entspricht dem
  \emph{Offline Longest Common Subsequence}-Problem auf $n$ Vektoren der Dimension $n$,
  das unter OV nicht in $O(n^{2-\varepsilon} \cdot n)$ lösbar ist
  \cite{abboud2014consequences}. Damit gilt $T_{\text{Subgraph}} \geq \Omega(n^{2-\varepsilon})$.
\end{proof}

\subsubsection{Unbedingte untere Schranke: Kommunikationskomplexität}

Für eine unbedingte (hypothesenfreie) $\Omega(n^2)$-Schranke kann ein
Kommunikationskomplexitäts-Argument verwendet werden \cite{kushilevitz1997communication}:
Alice hält Matrix $A$, Bob hält Matrix $B$. Um $G \subseteq G'$ zu entscheiden, muss
mindestens $\Omega(n^2)$ Information ausgetauscht werden (da die Eingabe $2n^2$ Bits
umfasst und die Ausgabe von allen Bits abhängen kann). Dies beweist ohne jede Hypothese:
\[
  T_{\text{Subgraph}}(n) \geq \Omega(n^2) \quad \text{(unbedingt)}.
\]

\subsection{Erweiterung auf gerichtete azyklische Multigraphen und Hypergraphen}

In modernen VLSI-Designs treten Strukturen auf, die über einfache Digraphen
hinausgehen:

\subsubsection{Multigraphen für mehrfach verbundene Netze}

In einem Routing-Graphen können zwei Zellen durch mehrere Verdrahtungsoptionen
verbunden sein (z.\,B.\ Metall-1, Metall-2, Metall-3 in einem Dreischicht-Prozess).
Ein \emph{Multigraph} erlaubt Mehrfachkanten zwischen denselben Knoten:
\[
  E \subseteq V \times V \times L, \quad L = \{\text{met1}, \text{met2}, \text{met3}\}.
\]
Die Adjazenzmatrix wird zur Adjazenzmatrix-Familie $\{A^{(l)}\}_{l \in L}$, eine für
jede Schicht. Der Subgraph Algorithmus lässt sich direkt auf Multigraphen erweitern,
indem die Signatur um Schicht-Terme ergänzt wird:
\[
  \sigma_j^{(l)} = \sum_{i=0}^{n-1} A_{ij}^{(l)} \cdot 2^i + j \cdot 2^n + \text{hash}(l) \cdot 2^{2n}.
\]
Für $|L|$ Schichten steigt die Laufzeit auf $O(|L| \cdot n^3)$ -- linear in der Anzahl
der Schichten.

\subsubsection{Hypergraphen für Multi-Fanout-Netze}

In Logiknetzen kann ein Gate-Ausgang an beliebig viele Gate-Eingänge verbunden sein
(Fanout $> 2$). Dies wird durch Hyperkanten modelliert. Ein
\emph{gerichteter Hypergraph} $H = (V, \mathcal{E})$ hat Hyperkanten
$e \in \mathcal{E}$ mit $e \subseteq V \times 2^V$. Die Adjazenzmatrix wird zur
\emph{Inzidenzmatrix} $B \in \{0,1\}^{n \times m}$ (Knoten $\times$ Hyperkanten).

Die Erweiterung des Signaturkonzepts auf Inzidenzmatrizen ist nicht-trivial, da
die LCS-basierte Rotationsprüfung auf rechteckigen Matrizen der Form $n \times m$
(mit $m \neq n$) operate. Erste Vorstudien legen nahe, dass für Hypergraphen mit
maximalem Fanout $f$ eine Laufzeit von $O(f \cdot n^3)$ erreichbar ist.

\subsection{Automatische Verifikation digitaler Systeme}

\subsubsection{Hardware-Model-Checking mit Subgraph-Unterstützung}

Model Checking \cite{clarke1986automatic} prüft, ob ein endlicher Automat $M$ eine
temporallogische Spezifikation $\varphi$ erfüllt. Ein zentrales Teilproblem ist die
\emph{Bisimulation}: Zwei Automaten $M_1$ und $M_2$ sind bisimulationsäquivalent genau
dann, wenn jeder Übergang von $M_1$ durch einen entsprechenden Übergang von $M_2$
simuliert werden kann und umgekehrt.

Bisimulation ist stärker als Subgraph-Inklusion (sie erfordert symmetrisches Matching),
aber der Subgraph Algorithmus kann als effizienter Vorfilter in einem zweistufigen
Verfahren dienen:
\begin{enumerate}
  \item \textbf{Vorstufe:} Prüfe $G_{M_1} \subseteq G_{M_2}$ und
        $G_{M_2} \subseteq G_{M_1}$ mit dem Subgraph Algorithmus in $O(n^3)$.
        Falls eines scheitert, ist Bisimulation ausgeschlossen.
  \item \textbf{Nachstufe:} Führe exaktes Bisimulations-Checking durch (z.\,B.
        Hopcroft-Karp \cite{hopcroft1971} in $O(n \log n)$).
\end{enumerate}
Dieses zweistufige Vorgehen eliminiert in der Praxis den teuren zweiten Schritt für
alle Nicht-Bisimulations-Paare und reduziert die durchschnittliche Laufzeit
signifikant.

\subsubsection{Property-Checking in Verilog-Designs}

Moderne Verilog-Syntheseflüsse beinhalten Assertion-Checking (SVA, SystemVerilog
Assertions). Eine Assertion $\forall t: P(s_t)$ über einen Signalpfad $s_t$ des
Designs lässt sich als \emph{Mustererkennung im STG} auffassen: Der STG des Designs
darf den Teilgraphen $G_{\neg P}$ (der die Verletzungszustände der negierten
Eigenschaft darstellt) nicht enthalten.

\begin{satz}[Assertion-Prüfung via Subgraph]\label{satz:assertion}
  Gegeben ein Verilog-Design mit STG $G_M$ ($n$ Zustände) und eine
  SVA-Assertion, die durch einen Constraintgraphen $G_{\neg P}$ ($m \leq n$ Knoten)
  repräsentiert wird. Die Assertion wird verletzt, genau wenn $G_{\neg P} \subseteq G_M$.
  Dieser Test hat Laufzeit $O(n^3)$.
\end{satz}

\begin{proof}
  Die Korrektheit folgt direkt aus der Definition des Subgraph Algorithmus
  \cite{subgraph2026}: Falls $G_{\neg P} \subseteq G_M$, gibt es eine
  Zustandsfolge im Design, die die Assertion verletzt. Die Laufzeit $O(n^3)$
  folgt aus Satz~\ref{satz:stg_opt}.
\end{proof}

\subsection{Maschinelles Lernen und neuronale Netzarchitekturen}

\subsubsection{Graph Neural Networks für Subgraph-Erkennung}

Neuronale Graphennetzwerke (Graph Neural Networks, GNNs) \cite{schur2022gnn_survey}
können auf Graphen trainiert werden, um Subgraph-Inklusion approximativ zu entscheiden.
Dies ist relevant, wenn eine hohe Durchsatzrate wichtiger ist als exakte Korrektheit
(z.\,B.\ beim schnellen Screening von Millionen von Designvarianten).

Ein GNN der Tiefe $L$ mit versteckter Dimension $d$ hat Inferenzlaufzeit
$O(L \cdot (n \cdot d + |E| \cdot d^2))$ -- für $d = O(1)$ und $L = O(1)$ linear in $n$.
Der Subgraph Algorithmus hingegen ist exakt und garantiert korrekt.

\textbf{Hybridansatz:} Ein zweistufiges System mit GNN-Vorfilter ($O(n)$) und
nachgelagertem exaktem Subgraph Algorithmus ($O(n^3)$) nur für vielversprechende
Kandidaten könnte lineare durchschnittliche Laufzeit bei garantierter Korrektheit
erreichen.

\subsubsection{Lernbasierte Signaturverbesserung}

Die Signatur-Funktion $\sigma_j = \sum_i A_{ij} \cdot 2^i + j \cdot 2^n$ ist in
\cite{subgraph2026} analytisch hergeleitet. Es wäre interessant zu untersuchen, ob
eine \emph{gelernte Einbettung} $\phi: \{0,1\}^n \to \mathbb{R}^k$ mit $k \ll n$
die Rotationsprüfung beschleunigen kann, falls Kollisionen (falsche LCS-Treffer die
keine echten Subgraphen sind) toleriert werden.

\subsection{Anwendungen in neuen Halbleiterprozessen}

\subsubsection{3D-ICs und Die-Stacking}

In dreidimensionalen integrierten Schaltkreisen (3D-ICs) werden mehrere Dies übereinander
gestapelt und mit Through-Silicon Vias (TSVs) verbunden. Der Verbindungsgraph ist kein
planarer Graph mehr, sondern ein Mehrebenengraph $G_{3D} = (V, E_{\text{intra}} \cup E_{\text{inter}})$,
wobei $E_{\text{intra}}$ Verbindungen innerhalb einer Schicht und $E_{\text{inter}}$
Die-übergreifende TSV-Verbindungen bezeichnen. Der Subgraph Algorithmus kann auf
jede Schicht separat angewendet werden ($O(n_i^3)$ je Schicht $i$) und die TSV-Struktur
als Kantenannotation einbeziehen ($O(k \cdot n^3)$ für $k$ Schichten mit je $n$ Knoten).

\subsubsection{Quantenschaltkreise}

Quantenschaltkreise werden als gerichtete Graphen über einem komplexen Alphabet
(unitäre Gatter) repräsentiert. Obwohl die Adjazenzmatrix über $\mathbb{C}$ definiert
ist, kann der Subgraph Algorithmus mit komplexen Signaturen angepasst werden:
\[
  \sigma_j^{\mathbb{C}} = \sum_{i=0}^{n-1} |U_{ij}|^2 \cdot p_i + j \cdot 2^n,
\]
wobei $p_i$ Primzahlen sind (zur Eindeutigkeit). Gleichheit zweier Signaturen ist dann
ein notwendiges (aber nicht hinreichendes) Kriterium für unitäre Äquivalenz zweier
Quantengatter-Anordnungen. Der Algorithmus dient als effizienter Vorfilter für
Quantenschaltungsoptimierung \cite{nielsen2000quantum}.

\subsection{Weiterführende Perspektiven für Zustandskodierung}

Die Ergebnisse aus Kapitel~3 eröffnen mehrere konkrete Forschungsrichtungen:

\subsubsection{Lernbasierte Kodierungsoptimierung}

Das Problem der optimalen Zustandskodierung ist NP-hart (Satz~\ref{satz:kodzkomp}).
Der Subgraph Algorithmus reduziert den Suchraum von $m!$ auf $m!/|\text{Aut}(G_M)|$.
Eine natürliche Erweiterung ist ein Machine-Learning-Ansatz: Trainiere ein neuronales
Netzwerk, das aus den STG-Eigenschaften direkt eine hochqualitative Kodierung
$\kappa$ vorhersagt, ohne den gesamten Suchraum zu durchsuchen. Der Subgraph
Algorithmus dient dann als effizienter Validator ($O(m^3)$) im Trainingsloop.

\textbf{Offene Frage:} Gibt es eine Klasse von STGs (z.\,B. reguläre Graphen), für
die das Zustandskodierungsproblem polynomial lösbar ist?

\subsubsection{Mehrziel-Kodierungsoptimierung}

In der Praxis sind Gatterkomplexität, Taktfrequenz und Verlustleistung konkurrierende
Ziele. Wir definieren:
\[
  \kappa^* = \arg\min_\kappa \bigl(\alpha \cdot W(\kappa) + \beta \cdot P(\kappa)
  + \gamma \cdot D(\kappa)\bigr),
\]
wobei $W$ das Kodierungsgewicht (Satz~\ref{def:hamming}), $P$ die Verlustleistung
(proportional zur Schaltaktivität) und $D$ die kritische Pfadverzögerung bezeichnen.
Der Subgraph Algorithmus kann alle drei Terme durch Matchingoperationen auf dem
kodierten STG berechnen. Die Pareto-Optimierung dieses Mehrzielproblems bleibt offen.

\subsection{Perspektiven für formale Äquivalenzprüfung}

\subsubsection{Inkrementelle Äquivalenzprüfung bei RTL-Änderungen}

In modernen Entwurfsumgebungen werden RTL-Beschreibungen kontinuierlich geändert
(z.\,B.\ agile Hardware-Entwicklung). Nach jeder Änderung muss die Äquivalenz zur
Vorversion geprüft werden. Der Subgraph Algorithmus eignet sich besonders gut für
\emph{inkrementelle} FEC:

\begin{satz}[Inkrementelle Äquivalenzprüfung]\label{satz:inkrementell_fec}
  Wenn eine RTL-Änderung lokal $k$ Knoten des Netzlistengraphen betrifft
  ($k \ll n$), kann die inkrementelle Subgraph-basierte Äquivalenzprüfung in Zeit
  $O(k \cdot n^2)$ durchgeführt werden, anstatt in $O(n^3)$ für den vollständigen
  Neutest.
\end{satz}

\begin{proof}
  Nur die Signaturen der $k$ geänderten Knoten müssen aktualisiert werden: $O(k \cdot n)$.
  Die LCS-Berechnung zwischen den $k$ neuen und allen $n$ alten Signaturen kostet
  $O(k \cdot n^2)$. Gesamtlaufzeit: $O(k \cdot n^2) \ll O(n^3)$ für $k \ll n$.
\end{proof}

\subsubsection{Integration in ISO-26262-konforme Entwurfsflüsse}

Die formale Äquivalenzprüfung ist in funktional-sicherheitskritischen Entwurfsflüssen
nach ISO~26262 (Automotive) und DO-178C (Avionics) verpflichtend. Der Subgraph
Algorithmus bietet gegenüber proprietären Werkzeugen den Vorteil des
\emph{beweisbaren Laufzeitbudgets}: $T_{\text{FEC}} = O(n^3)$ ist formal nachgewiesen,
nicht nur empirisch beobachtet. Dies ist ein entscheidender Vorteil für das Safety Case
Assessment, da Werkzeugqualifizierung ($\text{TQL-1}$ nach DO-330) einen
dokumentierten, verifizierten Algorithmus erfordert.

\textbf{Konkrete nächste Schritte:}
\begin{enumerate}
  \item Implementierung des zweistufigen FEC-Algorithmus als quelloffener
        Yosys-Pass \cite{wolf2013yosys} in C++.
  \item Benchmarking auf dem ITC-99-Testsatz (sequentielle Benchmark-Schaltkreise).
  \item Formale Verifikation der Implementierung mit einem Proof-Assistant (Coq, Lean).
\end{enumerate}

\subsection{Weiterführende Perspektiven für ATPG}

\subsubsection{ATPG für Mehrfachfehler}

Das klassische Stuck-at-Fehlermodell betrachtet Einzelfehler. In modernen
Nanometer-Prozessen treten jedoch gehäuft Mehrfachfehler auf (Multiple Simultaneous
Faults, MSF). Ein MSF ändert mehrere Kanten im Netzlistengraph gleichzeitig.

\begin{definition}[k-facher Stuck-at-Fehler]\label{def:kfach}
  Ein $k$-facher Stuck-at-Fehler $F_k = \{f_1, \ldots, f_k\}$ erzeugt eine
  Adjazenzmatrixänderung:
  \[
    A_{F_k} = A_N + \sum_{l=1}^{k} \varepsilon_l \cdot \Delta_{f_l},
    \quad \varepsilon_l \in \{-1, +1\},
  \]
  wobei $\|\Delta_{f_l}\|_1 = 1$ für jeden Einzelfehler.
\end{definition}

\begin{satz}[Subgraph Algorithmus für k-fache Fehler]\label{satz:kfachfehler}
  Ein $k$-facher Stuck-at-Fehler $F_k$ ist strukturell detektierbar genau dann, wenn
  $G_{F_k} \not\subseteq G_N$. Der Subgraph Algorithmus entscheidet dies in
  $O(n^3)$, unabhängig von $k$.
\end{satz}

\begin{proof}
  Für beliebiges $k$ gilt: $A_{F_k}$ unterscheidet sich von $A_N$ in höchstens $k$
  Einträgen. Die Signatur $\Sigma^{A_{F_k}}$ unterscheidet sich von $\Sigma^{A_N}$
  in höchstens $k$ Spalten. Die modifizierte Signatursequenz ist in $O(k \cdot n)$
  inkrementell berechnet. Der LCS-Test bleibt $O(n^2)$ je Rotation und $O(n^3)$
  gesamt.
\end{proof}

\subsubsection{Delay-Fehlermodell und Path-Delay-Testing}

Das Stuck-at-Modell abstrahiert strukturelle Fehler; reale Fertigungsfehler äußern
sich oft als erhöhte Signalverzögerungen (Path Delay Faults). Ein Path-Delay-Fehler
$f_{\text{del}}$ erhöht das Kantengewicht $d(e)$ im Timing-DAG um $\Delta d > 0$.

\begin{satz}[Integriertes ATPG-STA via Subgraph Algorithmus]\label{satz:atpgsta}
  Das kombinierte Structural / Timing-ATPG (Stuck-at + Path-Delay) kann als
  Subgraph-Matching auf dem gewichteten Timing-DAG $G_T$ formuliert werden.
  Der Subgraph Algorithmus mit gewichteter Signatur
  $\sigma_j^{(w)} = \sum_i A_{ij} \cdot \lfloor d(i,j)/\delta \rfloor \cdot 2^i$
  (analog zu Definition~\ref{def:weighted_match} im Ausblick der Originalarbeit)
  löst das Problem in $O(n^3 \cdot W)$ mit Quantisierungsbreite $W$.
\end{satz}

\begin{proof}
  Ein Path-Delay-Fehler ändert ein Kantengewicht $d(e^*) \to d(e^*) + \Delta d$.
  In der quantisierten Signatur ändert sich der entsprechende Term:
  $\lfloor (d(e^*) + \Delta d)/\delta \rfloor \neq \lfloor d(e^*)/\delta \rfloor$
  falls $\Delta d \geq \delta$.
  Damit ist $\sigma_j^{(w, \text{fault})} \neq \sigma_j^{(w)}$ -- der Subgraph
  Algorithmus detektiert den Fehler. Laufzeit: Signaturberechnung mit $W$ Bits
  pro Gewichtseintrag kostet $O(n^2 W)$; LCS-Test kostet $O(n^2)$ je Rotation,
  Gesamtlaufzeit $O(n^3 \cdot W)$.
\end{proof}

\subsubsection{Diagnosefähigkeit: Fehlerlokalisation statt nur Detektion}

Über die Detektion hinaus ist die \emph{Fehlerlokalisation} (Fault Diagnosis) wichtig:
Welcher konkrete Fehler $f \in \mathcal{F}$ ist in einem defekten Chip aufgetreten?
Die Diagnose vergleicht gemessene Testantworten mit simulierten Fehlerantworten.

\begin{satz}[Fehlerlokalisation via Subgraph Algorithmus]\label{satz:fehlerlokal}
  Für einen defekten Chip mit beobachteter Ausgabesequenz $\mathbf{y}^{\text{obs}}$
  kann die Diagnoseliste $\mathcal{D} = \{f \in \mathcal{F} \mid N_f \text{ erklärt }
  \mathbf{y}^{\text{obs}}\}$ durch Subgraph-Matching aller Fehlergraphen $G_f$ gegen
  den ``Beobachtungsgraph'' $G_{\text{obs}}$ (der aus $\mathbf{y}^{\text{obs}}$ rekonstruiert
  wird) in Zeit $O(|\mathcal{F}| \cdot n^3 / n) = O(n^4)$ ermittelt werden; mit
  Signatur-Caching amortisiert auf $O(n^3)$.
\end{satz}

\begin{proof}
  Der Beobachtungsgraph $G_{\text{obs}}$ wird aus der gemessenen Ausgabe rekonstruiert
  (Modell-basierte Diagnose): $O(n)$. Für jeden Fehler $f \in \mathcal{F}$ ($|\mathcal{F}|
  = O(n^2)$):
  \begin{itemize}
    \item Prüfe inkrementell $G_f \subseteq G_{\text{obs}}$: $O(n^2)$ mit gecachten Signaturen.
  \end{itemize}
  Gesamt: $O(n^2) \cdot O(n^2) / n = O(n^3)$ amortisiert (Signatur-Caching wie in
  Satz~\ref{satz:atpg_subgraph}).
\end{proof}

\subsection{Formale Offene Probleme}

Wir schließen den Ausblick mit einer erweiterten Liste formaler offener Fragen, von
denen jede eine eigenständige Forschungsarbeit rechtfertigt:

\begin{enumerate}
  \item \textbf{Unbedingte $\Omega(n^3)$-Schranke:} Kann das zyklische Subgraph-Problem
        ohne Rückgriff auf Komplexitätshypothesen (SETH, OV) in $\Omega(n^3)$ bewiesen
        werden? Dies ist ein offenes Problem der Fine-Grained Complexity.

  \item \textbf{Sub-kubische Algorithmen unter Einschränkungen:} Für welche
        Graphklassen (z.\,B.\ planare Graphen, Graphen mit beschränkter Baumweite,
        DAGs mit kleinem Fanin) lässt sich das Subgraph-Problem in $O(n^{2.5})$ oder
        besser lösen?

  \item \textbf{Online-Variante:} Gegeben ein wachsender STG (neue Zustände kommen
        inkrementell hinzu). Kann die Subgraph-Prüfung in $O(n^2)$ amortisierter Zeit
        pro Einfügeoperation aufrechterhalten werden (analog zu dynamischen Graphalgorithmen)?

  \item \textbf{Parallelität und Raumkomplexität:} Kann der Subgraph Algorithmus
        in der PRAM-Klasse $\mathcal{NC}^2$ (polylogarithmische Zeit, polynomiell viele
        Prozessoren) oder $\mathcal{NC}^1$ realisiert werden?

  \item \textbf{Randomisierte Beschleunigung:} Gibt es einen Las-Vegas-Algorithmus, der
        das zyklische Subgraph-Problem im Erwartungswert in $O(n^{2.5})$ löst?

  \item \textbf{Erweiterung auf Graphen verschiedener Größe mit $n_A \neq n_B$:}
        Der aktuelle Algorithmus ist auf $n_A \leq n_B$ ausgelegt. Die allgemeine
        Version (verschiedene Knotenmengen, Einbettung durch beliebige -- nicht nur
        zyklische -- Permutation) ist NP-vollständig
        \cite{cook1971complexity}. Die Grenze zwischen polynomiellen und
        NP-schweren Varianten ist noch nicht vollständig charakterisiert.

  \item \textbf{Optimale Zustandskodierung in Polynomialzeit:} Gibt es eine Klasse
        von FSM-Graphen, für die das optimale Kodierungsproblem (minimales
        Kodierungsgewicht $W(\kappa)$) in Polynomialzeit lösbar ist? Erste Kandidaten
        sind reguläre Graphen und Graphen mit beschränkter Treewidth.

  \item \textbf{Vollständige sequentielle Äquivalenzprüfung ohne SAT-Backend:}
        Kann die sequentielle Äquivalenzprüfung (Bisimulation) rein durch den Subgraph
        Algorithmus entschieden werden, ohne Rückgriff auf den Hopcroft-Karp-Algorithmus?
        Hierfür müsste die Bisimulationsrelation direkt aus der Signaturstruktur
        ablesbar sein.

  \item \textbf{ATPG für k-fache Fehler mit minimaler Testmenge:} Das Minimum-Set-Cover-Problem
        für $k$-fache Stuck-at-Fehler ist für $k \geq 2$ noch schwerer als für
        $k = 1$. Gibt es approximationserhaltende Reduktionen, die zeigen, dass die
        $(1-1/e)$-Approximation aus Satz~\ref{satz:mintest} auch für $k \geq 2$ gilt?

  \item \textbf{Integriertes Synthesis-Verification Framework:} Kann ein einheitliches
        Werkzeug entwickelt werden, das Logik-Synthese (Anwendung IV), formale
        Äquivalenz\-prüfung (Anwendung III) und ATPG (Anwendung VII) in einem
        gemeinsamen Subgraph Algorithmus-Kern verbindet -- mit einer Gesamtlaufzeit
        $O(n^3)$ für den kompletten Synthesis-and-Verify-Flow?
\end{enumerate}

% ===========================================================================
\begin{thebibliography}{99}

\bibitem{subgraph2026}
S.~Epp,
\emph{Der Subgraph Algorithmus},
Verfügbar unter \url{https://gitlab.com/epp-group/subgraph}

\bibitem{ciletti2011}
M.~D.~Ciletti,
\emph{Advanced Digital Design with the Verilog HDL},
2.~Aufl., Prentice Hall, 2011.

\bibitem{abboud2015tight}
A.~Abboud, A.~Backurs, V.~V.~Williams,
\emph{Tight Hardness Results for LCS and Other Sequence Similarity Measures},
in Proc.~IEEE FOCS, 2015, S.~59--78.

\bibitem{brayton1987mis}
R.~K.~Brayton, R.~Rudell, A.~Sangiovanni-Vincentelli, A.~R.~Wang,
\emph{MIS: A Multiple-Level Logic Optimization System},
IEEE Trans. Comput.-Aided Des. Integr. Circuits Syst., Bd.~6, Nr.~6, S.~1062--1081, 1987.

\bibitem{cook1971complexity}
S.~A.~Cook,
\emph{The Complexity of Theorem Proving Procedures},
in Proc.~ACM STOC, 1971, S.~151--158.

\bibitem{cormen2022}
T.~H.~Cormen, C.~E.~Leiserson, R.~L.~Rivest, C.~Stein,
\emph{Introduction to Algorithms},
4.~Aufl., MIT Press, 2022.

\bibitem{devadas1994logic}
S.~Devadas, A.~Ghosh, K.~Keutzer,
\emph{Logic Synthesis},
McGraw-Hill, 1994.

\bibitem{mcmillan1993symbolic}
K.~L.~McMillan,
\emph{Symbolic Model Checking},
Kluwer Academic Publishers, 1993.

\bibitem{hopcroft1971}
J.~Hopcroft, R.~Karp,
\emph{A Linear Algorithm for Testing Equivalence of Finite Automata},
Technical Report 114, Cornell University, 1971.

\bibitem{kahng2022vlsi}
A.~B.~Kahng, J.~Lienig, I.~L.~Markov, J.~Hu,
\emph{VLSI Physical Design: From Graph Partitioning to Timing Closure},
2.~Aufl., Springer, 2022.

% ---------- 24 erweiterte Quellen (nicht im Text zitiert) ----------

\bibitem{garey1979computers}
M.~R.~Garey, D.~S.~Johnson,
\emph{Computers and Intractability: A Guide to the Theory of NP-Completeness},
W.~H.~Freeman, 1979.

\bibitem{ullmann1976algorithm}
J.~R.~Ullmann,
\emph{An Algorithm for Subgraph Isomorphism},
J.~ACM, Bd.~23, Nr.~1, S.~31--42, 1976.

\bibitem{cordella2004subgraph}
L.~P.~Cordella, P.~Foggia, C.~Sansone, M.~Vento,
\emph{A (Sub)Graph Isomorphism Algorithm for Matching Large Graphs},
IEEE Trans.~Pattern Anal.~Mach.~Intell., Bd.~26, Nr.~10, S.~1367--1372, 2004.

\bibitem{valiant1979universality}
L.~G.~Valiant,
\emph{Universality Considerations in VLSI Circuits},
IEEE Trans.~Comput., Bd.~30, Nr.~2, S.~135--140, 1981.

\bibitem{sentovich1992sis}
E.~M.~Sentovich, K.~J.~Singh, L.~Lavagno, C.~Moon, R.~Murgai,
A.~Saldanha, H.~Savoj, P.~R.~Stephan, R.~K.~Brayton, A.~Sangiovanni-Vincentelli,
\emph{SIS: A System for Sequential Circuit Synthesis},
Technical Report UCB/ERL M92/41, UC Berkeley, 1992.

\bibitem{brayton2010abc}
R.~Brayton, A.~Mishchenko,
\emph{ABC: An Academic Industrial-Strength Verification Tool},
in Proc.~CAV, Lecture Notes in Computer Science, Bd.~6174, S.~24--40, 2010.

\bibitem{wolf2013yosys}
C.~Wolf,
\emph{Yosys Open SYnthesis Suite},
Proc.~Austrochip, 2013.
\url{https://yosyshq.net}

\bibitem{brglez1985combinational}
F.~Brglez, H.~Fujiwara,
\emph{A Neutral Netlist of 10 Combinational Benchmark Circuits and a Target
Translator in Fortran},
in Proc.~IEEE ISCAS, S.~663--698, 1985.

\bibitem{jaja1992introduction}
J.~Jájá,
\emph{An Introduction to Parallel Algorithms},
Addison-Wesley, 1992.

\bibitem{allison1986lcs}
L.~Allison, T.~I.~Dix,
\emph{A Bit-String Longest-Common-Subsequence Algorithm},
Inform. Process. Lett., Bd.~23, Nr.~6, S.~305--310, 1986.

\bibitem{smith1981identification}
T.~F.~Smith, M.~S.~Waterman,
\emph{Identification of Common Molecular Subsequences},
J.~Molecular Biology, Bd.~147, Nr.~1, S.~195--197, 1981.

\bibitem{wagner1974string}
R.~A.~Wagner, M.~J.~Fischer,
\emph{The String-to-String Correction Problem},
J.~ACM, Bd.~21, Nr.~1, S.~168--173, 1974.

\bibitem{vassilevska2018hardness}
V.~Vassilevska Williams,
\emph{On Some Fine-Grained Questions in Algorithms and Complexity},
in Proc.~ICM, Bd.~3, S.~3431--3472, 2018.

\bibitem{abboud2014consequences}
A.~Abboud, V.~V.~Williams,
\emph{Popular Conjectures Imply Strong Lower Bounds for Dynamic Problems},
in Proc.~IEEE FOCS, S.~434--443, 2014.

\bibitem{kushilevitz1997communication}
E.~Kushilevitz, N.~Nisan,
\emph{Communication Complexity},
Cambridge University Press, 1997.

\bibitem{clarke1986automatic}
E.~M.~Clarke, E.~A.~Emerson, A.~P.~Sistla,
\emph{Automatic Verification of Finite-State Concurrent Systems Using
Temporal Logic Specifications},
ACM Trans.~Program. Lang.~Syst., Bd.~8, Nr.~2, S.~244--263, 1986.

\bibitem{synopsys_dc}
Synopsys,
\emph{Design Compiler User Guide},
Synopsys Inc., Version 2023.03, 2023.

\bibitem{schur2022gnn_survey}
J.~Zhou, G.~Cui, S.~Hu, Z.~Zhang, C.~Yang, Z.~Liu, L.~Wang, C.~Li, M.~Sun,
\emph{Graph Neural Networks: A Review of Methods and Applications},
AI Open, Bd.~1, S.~57--81, 2020.

\bibitem{nielsen2000quantum}
M.~A.~Nielsen, I.~L.~Chuang,
\emph{Quantum Computation and Quantum Information},
Cambridge University Press, 2000.

\bibitem{lengauer1990combinatorial}
T.~Lengauer,
\emph{Combinatorial Algorithms for Integrated Circuit Layout},
Wiley, 1990.

\bibitem{lee1959algorithm}
C.~Y.~Lee,
\emph{An Algorithm for Path Connections and Its Applications},
IRE Trans.~Electron. Comput., Bd.~EC-10, Nr.~3, S.~346--365, 1961.

\bibitem{karp1972reducibility}
R.~M.~Karp,
\emph{Reducibility Among Combinatorial Problems},
in Complexity of Computer Computations, Plenum Press, S.~85--104, 1972.

\bibitem{dijkstra1959note}
E.~W.~Dijkstra,
\emph{A Note on Two Problems in Connexion with Graphs},
Numerische Mathematik, Bd.~1, Nr.~1, S.~269--271, 1959.

\bibitem{tarjan1972depth}
R.~E.~Tarjan,
\emph{Depth-First Search and Linear Graph Algorithms},
SIAM J.~Comput., Bd.~1, Nr.~2, S.~146--160, 1972.

% ---------- 17 weitere Quellen ----------

\bibitem{leiserson1991retiming}
C.~E.~Leiserson, J.~B.~Saxe,
\emph{Retiming Synchronous Circuitry},
Algorithmica, Bd.~6, Nr.~1--6, S.~5--35, 1991.

\bibitem{micheli1994synthesis}
G.~De~Micheli,
\emph{Synthesis and Optimization of Digital Circuits},
McGraw-Hill, 1994.

\bibitem{hachtel1996logic}
G.~D.~Hachtel, F.~Somenzi,
\emph{Logic Synthesis and Verification Algorithms},
Kluwer Academic Publishers, 1996.

\bibitem{roth1966diagnosis}
J.~P.~Roth,
\emph{Diagnosis of Automata Failures: A Calculus and a Method},
IBM J.~Res.~Dev., Bd.~10, Nr.~4, S.~278--291, 1966.

\bibitem{bushnell2000essentials}
M.~L.~Bushnell, V.~D.~Agrawal,
\emph{Essentials of Electronic Testing for Digital, Memory and
Mixed-Signal VLSI Circuits},
Kluwer Academic Publishers, 2000.

\bibitem{bryant1986bdd}
R.~E.~Bryant,
\emph{Graph-Based Algorithms for Boolean Function Manipulation},
IEEE Trans.~Comput., Bd.~C-35, Nr.~8, S.~677--691, 1986.

\bibitem{biere2009handbook}
A.~Biere, M.~Heule, H.~van~Maaren, T.~Walsh (Hrsg.),
\emph{Handbook of Satisfiability},
IOS Press, 2009.

\bibitem{moore1956gedanken}
E.~F.~Moore,
\emph{Gedanken-Experiments on Sequential Machines},
in Automata Studies, Princeton University Press, S.~129--153, 1956.

\bibitem{mealy1955method}
G.~H.~Mealy,
\emph{A Method for Synthesizing Sequential Circuits},
Bell Syst.~Tech.~J., Bd.~34, Nr.~5, S.~1045--1079, 1955.

\bibitem{bohm1994minimization}
M.~Böhm, E.~Speckenmeyer,
\emph{A Fast Parallel SAT-Solver -- Efficient Workload BalancingI},
Ann.~Math.~Artif.~Intell., Bd.~24, Nr.~1--4, S.~231--248, 1998.

\bibitem{shannon1938symbolic}
C.~E.~Shannon,
\emph{A Symbolic Analysis of Relay and Switching Circuits},
Trans.~Am.~Inst.~Electr.~Eng., Bd.~57, Nr.~12, S.~713--723, 1938.

\bibitem{aho1974design}
A.~V.~Aho, J.~E.~Hopcroft, J.~D.~Ullman,
\emph{The Design and Analysis of Computer Algorithms},
Addison-Wesley, 1974.

\bibitem{tseitin1968complexity}
G.~S.~Tseitin,
\emph{On the Complexity of Derivation in Propositional Calculus},
in Studies in Constructive Mathematics and Mathematical Logic,
Consultants Bureau, S.~115--125, 1968.

\bibitem{nam2008logic}
G.-J.~Nam, J.~Cong (Hrsg.),
\emph{Modern Circuit Placement: Best Practices and Results},
Springer, 2007.

\bibitem{weste2011cmos}
N.~H.~E.~Weste, D.~M.~Harris,
\emph{CMOS VLSI Design: A Circuits and Systems Perspective},
4.~Aufl., Addison-Wesley, 2011.

\bibitem{papadimitriou1994complexity}
C.~H.~Papadimitriou,
\emph{Computational Complexity},
Addison-Wesley, 1994.

\bibitem{goldreich2008computational}
O.~Goldreich,
\emph{Computational Complexity: A Conceptual Perspective},
Cambridge University Press, 2008.

\bibitem{knuth2011art}
D.~E.~Knuth,
\emph{The Art of Computer Programming, Vol.~4A: Combinatorial Algorithms,
Part~1},
Addison-Wesley, 2011.

\end{thebibliography}

\end{document}
